\documentclass[11pt,a4paper]{article}
\usepackage[margin=24mm]{geometry}
\usepackage{fontspec}
\usepackage[english,russian]{babel}
\setmainfont[Path=paper/v2/fonts/,Extension=.ttf,UprightFont=*-Regular,BoldFont=*-Bold,ItalicFont=*-Italic,BoldItalicFont=*-BoldItalic]{NotoSerif}
\setmonofont[Scale=.85,Path=paper/v2/fonts/,Extension=.ttf,UprightFont=*,BoldFont=*-Bold,ItalicFont=*-Oblique,BoldItalicFont=*-BoldOblique]{DejaVuSansMono}
\usepackage{amsmath,amssymb,amsthm}
\usepackage{array,booktabs,tabularx,enumitem}
\usepackage[table]{xcolor}
\definecolor{ink}{HTML}{1E293B}
\definecolor{slate}{HTML}{475569}
\definecolor{rulegray}{HTML}{CBD5E1}
\definecolor{tint}{HTML}{F8FAFC}
\definecolor{accent}{HTML}{174B69}
\usepackage[unicode,colorlinks=true,linkcolor=accent,urlcolor=accent,citecolor=accent,bookmarksnumbered=true]{hyperref}
\hypersetup{pdftitle={110 OBSERVER. Наблюдаемость собственного прошлого},pdfauthor={Aleksandr Strand},pdfsubject={Точное декодирование, Rule 110 и эксперименты на OLMo}}
\setlength{\emergencystretch}{2.5em}
\renewcommand{\contentsname}{Содержание}
\renewcommand{\refname}{Литература}
\renewcommand{\proofname}{Доказательство}
\renewcommand{\tablename}{Таблица}
\setcounter{tocdepth}{1}
\setlist{itemsep=1pt,topsep=3pt}
\newtheorem{theorem}{Теорема}[section]
\newtheorem{lemma}[theorem]{Лемма}
\newtheorem{proposition}[theorem]{Предложение}
\newtheorem{conjecture}[theorem]{Гипотеза}
\theoremstyle{definition}
\newtheorem{definition}[theorem]{Определение}
\newtheorem{remark}[theorem]{Замечание}
\newcommand{\src}[2]{\href{https://github.com/Dnarts/rule110-dependency-fronts/blob/v2.0-part-III-dev/#1}{#2}}
\begin{document}
\sloppy
\begin{center}
{\small\bfseries\color{accent}110 OBSERVER}\\[7pt]
{\LARGE\bfseries Наблюдаемость собственного прошлого}\\[5pt]
{\large Цена сообщения, причинные свидетели и границы вывода}\\[10pt]
Aleksandr Strand \hspace{1em} $\cdot$ \hspace{1em} 20 ноября 2026
\end{center}
\begin{quote}\small
\textbf{Срез.} Точная геометрия Rule 110 на бесконечной прямой имеет формальное доказательство в Lean 4.35.0. Данные OLMo зафиксированы на восьми серединах dev после двухшагового варианта A (\texttt{PASS\_DEV}). Панели 9--20, 21--32 и 33--44 имеют статус \texttt{NOT\_RUN} в этом срезе. Мост к феноменальному опыту имеет статус \texttt{DECLARED\_NOT\_DERIVED}.
\end{quote}
\noindent\textbf{Аннотация.} Один и тот же наблюдаемый ответ может возникнуть при чтении сохранённой истории, по корреляции или при постоянной привычке получателя. Мы формулируем условия, при которых сообщение различает нужные истории, и проверяем, воспользовалась ли им система. В Rule 110 изолированный переворот клетки даёт возможность проверить происхождение сигнала, а возврат сигнала после слияния конфигураций --- его действие. В замороженной языковой модели отправитель сводит выбор и обратную связь к одному токену до предъявления новой строки. С этим токеном получатель решает 64 из 64 различающих условий dev, с фиксированным токеном --- 32 из 64. В заключение указываем, какие дополнительные предпосылки потребовались бы для вывода о субъективном опыте.
\tableofcontents
\clearpage
\section{Постановка, две ключевые сцены и обозначения}
\emph{Прежде чем судить, на что система способна, проверьте, чего ей не дали узнать.}
Когда динамическая система может восстановить факт о собственном прошлом?
Сколько битов нужно ей передать и каким вмешательством можно проверить, что она
воспользовалась ими? Зафиксируем, какие величины доступны получателю к моменту
запроса. Две конструкции Rule~110 показывают, почему это существенно.

\paragraph{Два датчика покрывают окно и упускают различие.}
Поставим датчики в позициях $S=(-1,1)$ и прочтём их на тактах 0 и 1.
Нужно предсказать значение клетки 0 на такте 2. Всё окно её зависимости
$[-2,2]$ покрыто конусами датчиков. В первом начальном состоянии единица стоит только в $-1$,
во втором --- в $-1$ и $0$.
В обоих случаях датчики дважды показывают одну и ту же пару $(1,0)$.
На такте 2 целевая клетка равна соответственно 0 и 1. Какой бы ответ ни дал
декодер по этому журналу, в одном из двух случаев он ошибётся
(§\ref{sec:observer-coverage}).

\paragraph{Один бит возвращает утраченное различие.}
Теперь цель --- прежнее значение $Y$ клетки 5. В каждой из двух пар историй
переворот только $Y$ меняет показание $Z$ клетки 3 на такте 2. Наблюдатель
сохраняет этот бит. На такте 4 основные конфигурации внутри каждой пары
уже совпадают на всей прямой. Запишем $Z$ в клетку $-1$. Её левый сосед
равен 1, правый хранит контекстный бит $K$. Следующее обновление даёт
$Z\oplus K=Y$. Если при том же состоянии решётки вернуть $1-Z$, ответ
обратится во всех четырёх историях (§\ref{sec:observer-return}).

В этой конструкции по каналу идёт сырое показание датчика. В двухшаговом
опыте OLMo отправитель сначала по прошлому выбору $h$ и обратной связи $e$
вычисляет исправленное правило $p=h\oplus e$. В новый эпизод попадает только
выбранный до предъявления строки $\mathbf{x}$ токен $m\in\{\mathrm{A},\mathrm{B}\}$.
Получатель заново обрабатывает его в свежем KV-кэше. Так измеряется перенос
уже выбранного правила между эпизодами. Внутренняя коррекция стратегии при
непрерывном сохранении состояния требует отдельной проверки (§\ref{sec:olmo}).

\begin{center}\small
\begin{tabularx}{\linewidth}{@{}lX@{}}\toprule
Обозначение & Смысл на протяжении монографии \\\midrule
$\omega,H,W,Z$ & Допустимая история, сведения получателя, доступ наблюдателя и доставленное сообщение.\\
$Y$ & Целевой бит прошлого. В более общей постановке T1--T3 его алфавит может быть конечным.\\
$t,i,j,s$ & Дискретное время и целочисленные координаты решётки. $e_j$ меняет только начальную клетку $j$.\\
$B_t,\widetilde B_t,S_t$ & Невозмущённая и возмущённая конфигурации и множество координат, в которых они различаются.\\
$G_t,D_t,I_t$ & Значение клетки в начале координат, существенные начальные координаты и теоретическое окно зависимости.\\
$h,e,p,m$ & Прошлый выбор OLMo, бит обратной связи, исправленное правило и токен сообщения.\\
$\mathbf{x}=LwR$ & Новый восьмибитный вход. $w$ --- шестибитная середина, $L,R$ --- биты на концах.\\
$B,E(a,t)$ & Пакет мостовых постулатов и предикат опыта экземпляра $a$ в момент $t$.\\
\bottomrule\end{tabularx}
\end{center}

Формула указывает, что следует из заданных предпосылок (\texttt{COMPUTED}). Для вывода о данной машине нужен исполнимый свидетель или отчёт измерения (\texttt{MEASURED}). Статус \texttt{UNBOUND} означает отсутствие требуемой привязки, а \texttt{NOT\_RUN} --- отсутствие запуска в обозначенном срезе. Слово «точно» далее относится либо ко всем элементам заданной математической области, либо к явно указанной конечной панели. В §\ref{sec:observer-theory} сначала установим, при каких условиях точный декодер вообще существует.

\section{Доступ к прошлому: наблюдение, сообщение и память}
\label{sec:observer-theory}

Пусть $U\ne\varnothing$ --- множество допустимых историй одного устройства при
фиксированных границах системы и протоколе. Целевой бит $Y(u)\in\{0,1\}$
описывает прошлое устройства в истории $u$. Через $H(u)$ обозначим сведения,
доступные получателю к моменту решения, через $W(u)$ --- сведения отправителя.
Пришедшее сообщение обозначим $Z(u)$. В $H$ можно включать только те прежние
показания, которые доступны получателю по его интерфейсу или сохранены в его
памяти. Два протокола с разным доступом к журналу представляют разные задачи.

\subsection{T1. Критерий точного декодирования}
\label{thm:observer-t1}

Для произвольной карты наблюдения $O:U\to\mathcal O$ точное детерминированное
восстановление $Y$ на всех допустимых историях возможно тогда и только тогда,
когда
\begin{equation}
 \forall u,v\in U\qquad O(u)=O(v)\ \Longrightarrow\ Y(u)=Y(v).
 \label{eq:observer-fibres}
\end{equation}
\emph{Доказательство.} Если $Y=d\circ O$, применение одной функции $d$ к равным
аргументам даёт равные ответы. Обратно, для каждого $o\in O(U)$ множество
$\{Y(u):O(u)=o\}$ непусто и по условию одноэлементно. Определим $d(o)$ его
единственным элементом. Значения $d$ вне $O(U)$ произвольны. Это доказывает
существование функции. На конечной области её можно реализовать таблицей при
наличии памяти и адресации соответствующей ёмкости.

Критерий применим отдельно к $O=H$ и $O=(H,Z)$. Чтобы установить выигрыш от
сообщения в смысле безошибочного доступа, достаточно предъявить пару с общим
$H$ и разными $Y$, а затем дать точный декодер по $(H,Z)$ для \emph{всего}
класса историй. При заданном семействе наблюдений $O_t$ момент первого такого
доступа есть наименьшее $t$, при котором выполняется
\eqref{eq:observer-fibres}, если оно существует. В системе со стиранием
выполнение критерия в этот момент оставляет открытым его выполнение позднее.

\subsection{T2. Цена одного сообщения}
\label{thm:observer-t2}

Пусть $U$ конечно. Отправитель посылает одно слово $c(W)$ по безошибочному
каналу фиксированной длины. Моменты отправки и получения заданы заранее.
Построим неориентированный граф $G$ с вершинами $W(U)$. Между $w,w'$ проведём
ребро, если существуют $u,v\in U$ такие, что
\[
 W(u)=w,\qquad W(v)=w',\qquad H(u)=H(v),\qquad Y(u)\ne Y(v).
\]
Допускается $w=w'$. При наличии петли никакого достаточного сообщения такого
вида нет. В отсутствие петель минимальный размер алфавита и минимальное
число передаваемых двоичных разрядов равны соответственно
\[
 M_{\min}=\chi(G),\qquad b_{\min}=\lceil\log_2\chi(G)\rceil .
\]
\emph{Доказательство.} Для двух концов ребра совпадение сообщений породило бы
одинаковую пару $(H,Z)$ с различными целями, что противоречит T1. Следовательно,
значения всякого точного кода дают правильную раскраску. Петля потребовала бы
разных цветов у одной вершины. Обратно, при любой правильной раскраске две
истории с одинаковым $H$ и разными целями получают разные цвета. Если
отправитель передаёт цвет, T1 строит точный декодер. Для $M$ цветов требуется
и достаточно $\lceil\log_2 M\rceil$ битов фиксированной длины, включая
случай $M=1$ и пустого сообщения.

Нижняя граница $\max_{h\in H(U)}|\{Y(u):H(u)=h\}|$ достигается, если
отправитель знает и $H$, и $Y$. Он может локально пронумеровать возможные
цели для каждого $h$. При ограниченном доступе к истории эта оценка бывает
строго меньше $M_{\min}$. Пусть у получателя записана одна из упорядоченных
пар $ab,ac,bc$, отправитель знает одну выбранную букву этой пары, а целевой
бит указывает прежний выбор первой или второй позиции. Каждая пара требует
двух разных цветов для своих букв. Три буквы образуют треугольник, поэтому
здесь нужны три символа сообщения при двух возможных целевых битах.
Связь одношагового безошибочного кодирования со сторонней информацией и
раскраской графа восходит к
\href{https://doi.org/10.1109/TIT.1976.1055607}{работе Витсенхаузена};
функциональную постановку со сторонней информацией рассматривают также
\href{https://doi.org/10.3390/e26040338}{Шарпене, Ле Трёст и Руми}.

\subsection{T3. Ошибки канала и цена адресуемого запроса}
\label{thm:observer-t3}

Пусть в истории $u$ искажённый сигнал может принять любое значение из
непустого множества $\mathcal R(u)$. Восстановление при \emph{каждом}
допустимом искажении возможно ровно тогда, когда
\[
 H(u)=H(v),\quad Y(u)\ne Y(v)
 \quad\Longrightarrow\quad \mathcal R(u)\cap\mathcal R(v)=\varnothing .
\]
Общий символ из пересечения требует двух разных ответов на одно наблюдение.
Если все указанные пересечения пусты, каждому достижимому $(h,r)$
соответствует единственное значение цели. Это доказательство T1 на расширенных
историях с учётом исхода канала.

Для конечного совместного распределения $(H,R,Y)$ и штрафа за неверный бит
минимальная вероятность ошибки имеет другой смысл и равна
\[
 P_e^*=1-\sum_{h,r}\max_{y\in\{0,1\}}\Pr(H=h,R=r,Y=y).
\]
Действительно, детерминированный ответ $d(h,r)$ даёт вероятность успеха
$\sum_{h,r}\Pr(H=h,R=r,Y=d(h,r))$. Каждый член максимизируется независимо.
Например, если $H,Z$ --- равновероятные независимые биты,
$Y=H\oplus Z$ и $R=Z\oplus N$ с независимым $N\sim\mathrm{Bernoulli}(p)$,
$0\le p\le 1/2$, то оптимальный ответ $H\oplus R$ ошибается с вероятностью
$p$. При $p>0$ точная гарантия для всех искажений теряется.

Преимущество интерактивности зависит от того, какие сообщения разрешены.
Пусть $A_0,A_1$ --- два ранее записанных свободных бита, а локальный адрес
получателя $K$ определяет требуемый бит прошлого $Y=A_K$. Все восемь троек
$(K,A_0,A_1)$ допустимы. Достаточно одного бита запроса $K$ и одного
обратного бита $A_K$. Если обратный бит обязан быть фиксированной функцией
$c(A_0,A_1)$ независимо от адреса, то по одному его значению можно было бы
запустить декодер при $K=0$ и при $K=1$ и восстановить сразу пару
$(A_0,A_1)$. Четыре различные пары требуют не менее четырёх сообщений.
Таким образом, отправка адреса уменьшает необходимый обратный алфавит при
заданном устройстве обмена. Два обратных бита уже позволяют передать всю
запись без запроса. В рамках одного неизменного запроса заранее
воспроизведённый верный ответ с тем же содержанием и временем доставки
приводит к тому же исходу у детерминированного получателя.

\subsection{T4. Собственный отчёт и сохранённое состояние}
\label{thm:observer-t4}

Рассмотрим два устройства, которые вычисляют один и тот же точный
$d(H,Z)=Y$ и публикуют его. Первое сохраняет бит $M=Y$. Второе стирает
свидетельство и попадает в одно и то же доступное состояние хотя бы для
двух историй с разными $Y$. Затем прежний публичный отчёт и канал наблюдателя
закрываются, обоим устройствам предъявляется новый бит $C$ и запрашивается
действие $Y\oplus C$. Первое устройство вычисляет $M\oplus C$. Для второго
зафиксируем одно и то же $C$ на указанной паре историй. Полные доступные
входы равны, а требуемые биты противоположны. По T1 у второго устройства
точного декодера нет. Так две реализации с одинаковым немедленным отчётом
различаются отложенным действием.

Вмешательство в $Z$ при фиксированном раннем состоянии проверяет зависимость
действия от доставленного сигнала. Для проверки переноса результата через
эпизод требуется второе вмешательство: подмена или стирание $M$ при одном
и том же новом запросе и исключённом доступе к прежнему отчёту. Это два
разных контрфактических сравнения. На приведённой задаче одного внутреннего
регистра достаточно для успешного отложенного действия.

\section{Точная геометрия существенных координат Rule 110}\label{sec:geometry}
\subsection{Наблюдаемая и однобитовые свидетельства}
Фиксируем одну целевую клетку в начале координат и горизонт $t$. Входная координата существенна, если найдётся начальная конфигурация, в которой изменение одного этого бита меняет ответ. Фон может зависеть как от координаты, так и от горизонта. Это различие кванторов отделяет предельную чувствительность от скорости конкретного глайдера. Для Rule~110 известны конструкция универсального вычисления \cite{cook2004} и классификация движущихся структур \cite{wolfram1986,martinez2006}. Связь наблюдений с предсказательными классами и ограничениями вывода обсуждается в \cite{shalizi2001,wolpert2008}.

\begin{definition}[Rule 110 и целевая наблюдаемая]
Положим $X = \{0, 1\}^{\mathbb{Z}}$ и зададим локальное правило над полем $\mathbb{F}_2$ алгебраической нормальной формой
\begin{equation}\label{eq:rule110}
f(l, c, r) = c \oplus r \oplus cr \oplus lcr.
\end{equation}
Синхронный оператор эволюции $F : X \to X$ определяется равенством $(F(u))_i = f(u_{i-1}, u_i, u_{i+1})$. Наблюдаемая в начале координат в момент времени $t \in \mathbb{N}$ есть $G_t(u) = (F^t(u))_0$.
\end{definition}

\begin{definition}[Предсказательная факторизация и существенные переменные]
Отображение наблюдения $O : X \to \mathcal{O}$ называется достаточным для $G_t$, если существует отображение $P : \mathcal{O} \to \{0, 1\}$ такое, что $G_t = P \circ O$. Для $j \in \mathbb{Z}$ определим однобитовое возмущение $(u \oplus e_j)_i = u_i \oplus \mathbf{1}_{i=j}$. Множество существенных начальных координат определяется как
\begin{equation}\label{eq:essential_set}
D_t = \{ j \in \mathbb{Z} : \exists u \in X,\; G_t(u) \ne G_t(u \oplus e_j) \}.
\end{equation}
\end{definition}

Поскольку $G_t$ зависит лишь от конечного числа координат, $D_t$ является единственным минимальным по включению множеством входных координат, достаточным для вычисления $G_t$, а любой достаточный связный интервал наблюдения обязан содержать целочисленную выпуклую оболочку множества $D_t$.

Для работы со свидетельствами на периодических фонах удобно фиксировать связь между входной координатой $j$ и физической координатой $s$ дефекта на решётке. Для периодического слова $w$, пространственной фазы $\phi$ и соответствующей бесконечной конфигурации $B_0^{w,\phi}$ положим
\begin{equation}\label{eq:diff_support}
B_t^{w,\phi} = F^t(B_0^{w,\phi}), \qquad
\widetilde B_t^{w,\phi} = F^t(B_0^{w,\phi} \oplus e_0), \qquad
S_t^{w,\phi} = \{ s \in \mathbb{Z} : B_t^{w,\phi}(s) \ne \widetilde B_t^{w,\phi}(s) \}.
\end{equation}
В силу сдвиговой эквивариантности $F \circ \sigma^d = \sigma^d \circ F$, где $(\sigma^d u)_i = u_{i-d}$, любая точка $s \in S_t^{w,\phi}$ порождает свидетельство $j = -s \in D_t$ после сдвига начальной пары на $-s$. Следовательно, для любого семейства фонов и фаз $\mathcal{W}$ выполняется включение
\begin{equation}\label{eq:coverage_bridge}
C_t(\mathcal{W}) = \bigcup_{(w,\phi) \in \mathcal{W}} \{ -s : s \in S_t^{w,\phi} \} \subseteq D_t.
\end{equation}
При этом левая граница входного окна $j = -L(t)$ соответствует правому физическому фронту дефекта $s = L(t)$, а правая граница входного окна $j = t$ — левому физическому фронту $s = -t$.

\subsection{Точное окно зависимости на бесконечной решётке}

Положим $L(t) = t - \lfloor t/3 \rfloor$ и $I_t = [-L(t), t] \cap \mathbb{Z}$.

\begin{lemma}[Трёхтактное правостороннее исключение]\label{lem:three_step}
Для любого семибитного слова $b = (b_0, \dots, b_6) \in \{0, 1\}^7$ обозначим через $Q(b)$ выход центрального узла после трёх синхронных шагов Rule 110. Тогда
\begin{equation}\label{eq:seven_bit}
Q(b_0, b_1, \dots, b_6) = Q(1 \oplus b_0, b_1, \dots, b_6).
\end{equation}
\end{lemma}

\begin{proof}
Из \eqref{eq:rule110} булева производная по левому аргументу равна
\begin{equation}\label{eq:bool_deriv}
f(1, c, r) \oplus f(0, c, r) = cr.
\end{equation}
Пусть $b_i(t)$ обозначает значение клетки $i$ в одной из рассматриваемых траекторий. Предположим, что две конфигурации строго совпадают правее узла $p$, и их крайнее правое различие максимально смещается вправо за три такта: $p \to p+1 \to p+2 \to p+3$. Первый сдвиг вправо требует, чтобы общие значения в $p+1$ и $p+2$ при $t=0$ равнялись $1$; второй сдвиг требует значений $1$ в $p+2$ и $p+3$ при $t=1$. Поскольку $f(1, 1, z) = 1 \oplus z$, это вынуждает $b_{p+3}(0) = 0$. Далее, так как $f(1, 0, z) = z$, для поддержания фронта необходимо $b_{p+4}(0) = 1$. Но тогда $b_{p+4}(1) = f(0, 1, b_{p+5}(0)) = 1$, откуда $b_{p+3}(2) = f(1, 1, 1) = 0$, что противоречит третьему максимальному сдвигу вправо. Тождество \eqref{eq:seven_bit} также исчерпывающе проверено на всех $2^7 = 128$ входах в ядре Lean 4.
\end{proof}

\begin{theorem}[Универсальное включение множества зависимости]\label{thm:upper_bound}
Для каждого $t \in \mathbb{N}$ справедливо включение $D_t \subseteq I_t = [-(t - \lfloor t/3 \rfloor), t] \cap \mathbb{Z}$.
\end{theorem}

\begin{proof}
Если две конфигурации совпадают строго правее $p_0$, то за один шаг согласие сохраняется строго правее $p_0 + 1$, а по Лемме~\ref{lem:three_step} за три шага — строго правее $p_0 + 2$. Представляя $t = 3q + r$ с $q = \lfloor t/3 \rfloor$ и $r \in \{0, 1, 2\}$, получаем сохранение согласия правее $p_0 + 2q + r = p_0 + L(t)$. Следовательно, различие в точке $j < -L(t)$ не достигает начала координат к моменту $t$. При $j > t$ распространение блокируется стандартным световым конусом радиуса $1$.
\end{proof}

\begin{theorem}[Достижимость обеих границ]\label{thm:sharpness}
Для каждого $t \in \mathbb{N}$ обе крайние точки $-L(t)$ и $t$ принадлежат $D_t$. Минимальный достаточный связный интервал наблюдения равен $I_t$ и имеет ширину
\begin{equation}
m_{\min}(t) = 2t + 1 - \lfloor t/3 \rfloor, \qquad \lim_{t \to \infty} \frac{m_{\min}(t)}{t} = \frac{5}{3}.
\end{equation}
\end{theorem}

\begin{proof}
Для правой границы $j = t$ достаточно сравнить нулевую конфигурацию $0^{\mathbb{Z}}$ с одиночной единицей $e_t$: в силу $f(0, 0, 0) = 0$ и $f(0, 0, 1) = 1$ крайняя левая единица смещается ровно на одну клетку влево за каждый такт и достигает нуля за $t$ шагов.

Для левой границы $j = -L(t)$ рассмотрим фон периода четыре $B_0(i) = (1, 0, 1, 1)[i \bmod 4]$ и возмущение $\widetilde B_0 = B_0 \oplus e_2$. Прямая итерация Rule 110 на шести тактах даёт эволюцию носителя различий $S_t$, представленную в таблице~\ref{tab:period4_defect}.

\begin{table}[h]
\centering
\caption{Эволюция одиночного дефекта на фоне $(1011)$ и периодический сдвиг с шагом $(\Delta t, \Delta s) = (3, +2)$}
\label{tab:period4_defect}
\begin{tabular}{cccc}
\toprule
$t$ & Слово фона $B_t$ (период $4$) & Носитель различий $S_t$ & $\max S_t$ \\
\midrule
$0$ & $1011$ & $\{2\}$ & $2$ \\
$1$ & $1110$ & $\{1, 3\}$ & $3$ \\
$2$ & $1011$ & $\{1, 3, 4\}$ & $4$ \\
$3$ & $1110$ & $\{0, 1, 4\}$ & $4$ \\
$4$ & $1011$ & $\{-1, 0, 1, 3, 5\}$ & $5$ \\
$5$ & $1110$ & $\{3, 5, 6\}$ & $6$ \\
$6$ & $1011$ & $\{2, 3, 6\}$ & $6$ \\
\bottomrule
\end{tabular}
\end{table}

Из локальности эволюции и совпадения срезов при $t=3$ и $t=6$ со сдвигом на $2$ клетки вправо следуют точные тождества конфигураций $B_6 = \sigma^2 B_3$ и $\widetilde B_6 = \sigma^2 \widetilde B_3$. По сдвиговой эквивариантности $(B_{3+3k+r}, \widetilde B_{3+3k+r}) = \sigma^{2k}(B_{3+r}, \widetilde B_{3+r})$ для всех $k \ge 0$ и $r \in \{0, 1, 2\}$, откуда $\max S_t = 2 + L(t)$ для всех $t \ge 0$. Сдвиг начальной пары на $-[2 + L(t)]$ переносит начальное возмущение в $-L(t)$, а различие в момент $t$ — в $0$, что доказывает $-L(t) \in D_t$.
\end{proof}

\subsection{Конечные кольца и расхождение между ёмкостью и доступом}

На конечном периодическом кольце $\mathbb{Z}_N$ ($N \ge 3$) с пространством конфигураций $X_N = \{0, 1\}^{\mathbb{Z}_N}$ и оператором $F_N : X_N \to X_N$ обозначим через $\mathcal{R}_t^{(N)} = F_N^t(X_N)$ множество достижимых состояний, а через $C_t^{(N)} = \lceil \log_2 |\mathcal{R}_t^{(N)}| \rceil$ — битовую ёмкость состояния. Введём спаренную динамику $\Phi_N(u, v) = (F_N(u), F_N(v))$, начальные пары $E_j = \{ (u, u \oplus e_j) : u \in X_N \}$ и множество различий в начале координат $W_0 = \{ (u,v) : u_0 \ne v_0 \}$. Тогда
\begin{equation}
j \in D_t^{(N)} \iff \Phi_N^t(E_j) \cap W_0 \ne \emptyset.
\end{equation}

\begin{proposition}[Конечный сертификат персистентной существенности]\label{prop:ring_cert}
Пусть $\mu_N$ — максимальная длина предпериода конечной детерминированной системы $F_N$, а $P_N$ — наименьшее общее кратное длин её циклов. Тогда $D_{t+P_N}^{(N)} = D_t^{(N)}$ при всех $t \ge \mu_N$. Следовательно, для любого $T \in \mathbb{N}$ условие $\forall t \ge T,\; D_t^{(N)} = \mathbb{Z}_N$ эквивалентно проверке $\Phi_N^t(E_j) \cap W_0 \ne \emptyset$ для всех $j \in \mathbb{Z}_N$ и всех $T \le t < \max(T, \mu_N) + P_N$.
\end{proposition}

\begin{proof}
Каждое достижимое состояние $F_N^t(u)$ при $t\ge\mu_N$ лежит на цикле, длина которого делит $P_N$. Следовательно, отображения $F_N^{t+P_N}$ и $F_N^t$ совпадают на всём $X_N$. Для любого $j$ это сохраняет совпадение или различие во всех парах $(u,u\oplus e_j)$. С момента $\max(T,\mu_N)$ достаточно проверить один период. Если $T<\mu_N$, дополнительно проверяются все времена от $T$ до $\mu_N-1$.
\end{proof}

\begin{remark}[Компрессия состояний и чувствительность к входам]
Локальная чувствительность на бесконечной решётке не гарантирует персистентной существенности на конечном кольце: для Rule 110 при $N = 5$ все $32$ состояния схлопываются в $0^5$ к моменту $t = 5$, давая $D_t^{(5)} = \emptyset$ при $t \ge 5$. Напротив, полный перебор для Rule 110 на кольце $N = 20$ при $t = 121$ обнаруживает $|\mathcal{R}_{121}^{(20)}| = 2095$ достижимых состояний ($12$ бит ёмкости хранения), но при этом $D_{121}^{(20)} = \mathbb{Z}_{20}$ — все $20$ входных бит остаются существенными. Это разделяет четыре независимых аспекта задачи предсказания: \emph{причинность} (носитель $D_t$), \emph{доступ} (индивидуальное чтение координат против предобработки), \emph{ёмкость} ($C_t^{(N)}$) и \emph{вычислительную сложность}.
\end{remark}

\begin{theorem}[Сплошной носитель существенности]\label{thm:gap-free}
Для каждого $t \in \mathbb{N}$ выполняется равенство $D_t = I_t = [-(t - \lfloor t/3 \rfloor), t] \cap \mathbb{Z}$.
\end{theorem}

Первые конечные проверки охватили $0\le t\le2000$. Более поздний эксперимент с фиксированными фазами семи слов продолжил проверку до $t=5000$. Доказательство для всех времён завершается в подразделе~\ref{sec:five-bands}. Там используется новый конечный сертификат с границей $t=2484$.

\subsection{Односторонние автономные тайлы и две бесконечные приграничные полосы}

Из Леммы~\ref{lem:three_step} композицией вытекает, что область зависимости $3m$-тактного выхода в ячейке $s$ удовлетворяет включению $\mathrm{Dep}_{3m}(s) \subseteq [s - 2m, s + 3m]$ для любых конфигураций.

\begin{lemma}[Одностороннее полулучевое замыкание]\label{lem:one_sided}
Пусть $A_t = F^t(A_0)$, $p > 0$ и $a(t + p) = a(t) + \delta$. Если на базовых классах $t_0 \le t < t_0 + p$ выполнено $A_{t+p}(s + \delta) = A_t(s)$ для всех $s \ge a(t)$, и вычисление $A_{t+p}$ на полулуче $[a(t+p), \infty)$ зависит только от значений $A_t$ на $[a(t), \infty)$, то сдвиговое равенство $A_{t+p}(s + \delta) = A_t(s)$ сохраняется для всех $t \ge t_0$ и $s \ge a(t)$. Аналогичное утверждение верно для левых полулучей $s \le a(t)$.
\end{lemma}

Применяя Лемму~\ref{lem:one_sided} к правому физическому фронту $e(t) = L(t)$ с шагом $(\Delta t, \Delta s) = (30, +20)$ (где $e(t+30) = e(t) + 20$, а область зависимости нового тайла ширины $W$ лежит в $[e(t) - W + 1, e(t) + 50]$ и не затрагивает неизвестную внутреннюю область $s < e(t) - W + 1$) и аналогично к левому фронту $s = -t$ с шагом $(\Delta t, \Delta s) = (40,-40)$, получаем первую крупную бесконечную теорему.

\begin{theorem}[Две приграничные полосы]\label{thm:dual_strips}
Для каждого целого $t \ge 300$ выполнено
\begin{equation}\label{eq:dual_strips}
[-L(t), -L(t) + 105] \cap \mathbb{Z} \;\cup\; [t - 7, t] \cap \mathbb{Z} \;\subseteq\; D_t.
\end{equation}
\end{theorem}

Левая полоса ширины $106$ клеток сертифицируется $20$ фазами фонов $(0001111001)$, $(1011)$ и $(001011)$ с глубиной тайла $167$ клеток на $30$ базовых тактах $t = 300, \dots, 329$. Правая полоса ширины $8$ клеток сертифицируется шестью фазами с шагом $(\Delta t, \Delta s) = (40,-40)$ на базовых тактах $t = 300, \dots, 339$. Происхождение промежуточных состояний удостоверено по схеме \emph{proof-carrying computation} на битовых векторах \texttt{BitVec 1320} с проверкой каждого перехода $F_{1320}^{h_i}(z_i) = z_{i+1}$ ($h_i \in \{1, 10\}$) в ядре Lean 4.35.0.

\subsection{Баллистические лучи и границы локальных сертификатов}

Второй механизм неограниченного продолжения во времени основан на сохранении различия вдоль дискретных мировых линий.

\begin{lemma}[Сохранение правого полулуча]\label{lem:ballistic_ray}
Если для траектории $A$ в момент $T$ выполнено $\forall s \ge s,\; A_{T+3}(s + 2) = A_T(s)$, то $\forall q \ge 0,\; \forall s \ge s,\; A_{T+3q}(s + 2q) = A_T(s)$. В частности, если это условие выполнено одновременно для $B$ и $\widetilde B$ и $B_T(s) \ne \widetilde B_T(s)$, то $B_{T+3q}(s + 2q) \ne \widetilde B_{T+3q}(s + 2q)$ для всех $q \ge 0$.
\end{lemma}

Для двух арифметических серий якорей на фонах $(0000101111)$ и $(001011)$ при $0 \le n \le 32$:
\begin{equation}
T_A(n) = 594 + 882n,\quad s_A(n) = 256 + 378n; \qquad
T_B(n) = 884 + 882n,\quad s_B(n) = 380 + 378n,
\end{equation}
вычислительно сертифицированы $66$ начальных точек, порождающих бесконечные баллистические лучи существенности на постоянных вдоль каждого луча глубинах $d_A(n) = 140 + 210n$ и $d_B(n) = 210 + 210n$ от фронта $e(t) = L(t)$.

В последующих модулях \texttt{StaircaseA} и \texttt{StaircaseB} формулы обеих серий доказаны для произвольных $n,q\ge0$. Для этого три локально согласованные строки склеиваются в точную эволюцию. На передних периодических средах серий A и B обнаружено $45\,255$ покрытых классов из $66\,150$; с третьей серией --- $53\,550$. Число классов относится к таблице остатков $315\times210$. Применение таблицы в физической точке требует также проверки границ соответствующего клина.

Важную методологическую роль сыграли строгие отрицательные результаты, уточнившие границы применимости простых гипотез:
\begin{itemize}
    \item \textbf{Опровержение локально-тайловой макроиндукции по $n \mapsto n+1$:} при переходе на макрошаг $(\Delta T, \Delta s) = (882, +378)$ универсальная область зависимости составляет $\mathrm{Dep}_{882}(s) \subseteq [s - 588, s + 882]$ (прообраз $[s - 210, s + 1260]$). Прямыми контрпримерами в \texttt{check\_macro\_boundary.py} показано, что фиксация локальных тайлов радиуса $64$ (для $T=594$) или $186$ (для $T=884$) недостаточна: изменение клеток вне тайла меняет целевой бит через $882$ такта.
    \item \textbf{Неполнота наивных угловых секторов:} для фона $(0001111001)$ наблюдаемый до $t = 2000$ левый сектор со скоростями $[-0.66, -0.64]$ обнаруживает первый пропуск на более длинном горизонте.
    \item \textbf{Разделение периодичности и полноты покрытия:} глубокая сдвиговая рекуррентность ($261$ клетка для $(0000001001)$) не гарантирует сплошного покрытия без объединения нескольких фаз по всем классам вычетов.
\end{itemize}

\subsection{Конструктивное замыкание бесконечного левого клина: теорема \texttt{wedge\_all}}

Чтобы преодолеть несоизмеримость скоростей между левым краем светового конуса (скорость $-1$), фоновым кристаллом периода $50$ и фронтом распространения возмущения, вводится дискретная параметризация чётных тактов $n \in \mathbb{N}$:
\begin{equation}
\mathrm{tm}(n) = 2450 + 2n, \qquad a(n) = -350 - n.
\end{equation}
Здесь правая граница клина $a(n)$ смещается влево со скоростью $-1/2$ клетки за такт ($1$ клетка за $2$ такта). Для каждой из шести фаз $p \in \{\mathrm{p0}, \mathrm{p1}, \mathrm{p2}, \mathrm{p3}, \mathrm{p7}, \mathrm{p9}\}$ пишем $\widetilde B_p(t,s)=\widetilde B^p_t(s)$. Сопряжённый инвариант на слое $n$ объединяет совпадение с фоновым кристаллом $\mathrm{Cm}_p$ на интервале слева от $a(n)$ и принадлежность $16$-битного граничного буфера узлу конечного графа переходов:
\begin{align}
\mathrm{Medium}(p,n)\ \iff\ {}&
\forall s\in[-\mathrm{tm}(n)+1300,a(n))\cap\mathbb{Z},\notag\\
&\widetilde B_p(\mathrm{tm}(n),s)=\mathrm{Cm}_p(\mathrm{tm}(n),s), \\
\mathrm{Window}(p,n,i)\ \iff\ {}&
\forall k\in\{0,\dots,15\},\notag\\
&\widetilde B_p(\mathrm{tm}(n),a(n)+k)=\mathrm{testBit}(\mathrm{word}(p,n,i),k).
\end{align}

\begin{lemma}[16-битный детерминированный трансдьюсер и склейка со швом]\label{lem:transducer}
Для каждой фазы $p$ таблица \texttt{Crystal50Packed} задаёт $50$-периодический граф $36$-битных записей (младшие $16$ бит кодируют окно $\mathrm{word}(p, n, i)$, биты $16\dots 25$ и $26\dots 35$ — переходы при правом бите $r \in \{0, 1\}$).
\begin{enumerate}[label=(\alph*)]
    \item На шаге $n$ $4$-битный левый контекст $\mathrm{ctx}(p, n \bmod 50)$ в точках $a(n) - 4 + m$ ($m \in \{0, 1, 2, 3\}$) имеет глубину $t + s = 2096 + n + m$ и принадлежит уже доказанному интервалу $\mathrm{Medium}(p, n)$.
    \item Двукратное побитовое сжатие $\mathrm{shrink}(\mathrm{shrink}(\mathrm{input}(p, n, i, r)))$ вычисляет $17$ бит конфигурации в момент $\mathrm{tm}(n+1)$ на $[a(n)-2, a(n)+14]$: младший бит ($k=0$) даёт $\mathrm{Cm}_p(\mathrm{tm}(n+1), a(n) - 2)$, а старшие $16$ бит ($k=1\dots 16$) совпадают с узлом $z$ следующего слоя графа в позиции $a(n+1) = a(n) - 1$.
    \item Левый край полосы $-\mathrm{tm}(n+1) + 1300 \le s < -\mathrm{tm}(n) + 1302$ имеет инвариантную глубину $k = s + \mathrm{tm}(n+1) \in \{1300,1301,1302,1303\}$ и закрывается теоремой \texttt{seam} (\texttt{Crystal50Seam}) благодаря инвариантности $t + s$ при сдвиге $(t, s) \mapsto (t + 50, s - 50)$.
\end{enumerate}
\end{lemma}

\begin{theorem}[\texttt{Rails.Crystal50Wedge.wedge\_all}]\label{thm:wedge_all}
Для любого натурального $t \ge 2485$ и любого целого $s$, удовлетворяющего неравенствам
\begin{equation}\label{eq:wedge_region}
-t + 1302 \le s \le 59 - \lfloor t/2 \rfloor,
\end{equation}
выполняется предикат \texttt{EssentialAt}\,$t\,(-s)$. Верхней границы по времени нет.
\end{theorem}

\begin{proof}[Схема доказательства]
Положим $n = \lfloor (t - 2450) / 2 \rfloor$. При $t \ge 2485$ имеем $n \ge 17$, а правая граница $59 - \lfloor t/2 \rfloor = -1166 - n = a(n) - 816$. Для чётных тактов $t = \mathrm{tm}(n)$ равенство $\widetilde B_p(t, s) = \mathrm{Cm}_p(t, s)$ следует непосредственно из инварианта $\mathrm{Medium}(p, n)$ для всех шести фаз $p$. Для нечётных тактов $t = \mathrm{tm}(n) + 1$ одношаговый конус зависимости $[s - 1, s + 1]$ целиком лежит внутри $[-\mathrm{tm}(n) + 1300, a(n))$, и совпадение переносится леммой локальности \texttt{carry}. Применение моста \texttt{essential\_of\_cm} завершает доказательство.
\end{proof}

\subsection{Полное покрытие на бесконечной прямой}\label{sec:five-bands}
В координате входа $j=-s$ положим $b=\lfloor t/2\rfloor$. При $t\ge2485$ последовательно применяются следующие пять результатов. Обе границы в каждой строке включены.
\begin{center}
\begin{tabular}{lll}
\toprule
Начало & Конец & Основание \\
\midrule
$-L(t)$ & $-L(t)+105$ & Полулучевой сертификат ширины 106 \\
$-L(t)+106$ & $b-60$ & \texttt{EtherWedge.half\_sector} \\
$b-59$ & $t-1302$ & \texttt{Crystal50Wedge.wedge\_all} \\
$t-1301$ & $t-20$ & \texttt{LeftCrystal50.left\_strip} \\
$t-19$ & $t$ & \texttt{right\_boundary\_width20} \\
\bottomrule
\end{tabular}
\end{center}
Три фазы эфира периода 14 дают вторую строку уже при $t\ge210$. Последняя строка усиливает раннюю правую полосу ширины восемь. При $t=2485$ третья строка состоит из единственной клетки $j=1183$. Начало каждой следующей строки на единицу больше конца предыдущей. Лемма \texttt{Rails.late\_all} объединяет пять доказательств.

\paragraph{Конечный сертификат.}
Для $0\le t<2485$ берём все циклические фазы семи слов. Их длины и целые кодировки (младший бит соответствует первой клетке) равны
\[
(8,244),\ (10,632),\ (10,976),\ (10,576),\ (4,13),\ (6,52),\ (6,48).
\]
Фоны представлены векторами длины 5520; однобитный дефект внесён в позицию 2760. Периоды всех семи слов делят 5520. Для каждой строки времени строим побитовое объединение различий всех пар и проверяем, что оно содержит маску физического интервала $[-t,L(t)]$. Это $\sum_{t=0}^{2484}(2t+1-\lfloor t/3\rfloor)=5\,147\,263$ пары. При $s$ внутри интервала выполнено $-5520<s-t$ и $s+t<5520$. Поэтому лемма \texttt{local\_obs} переносит результат на бесконечную прямую с единственным переворотом.

Начальная строка вычисляется из семи слов. Генератор предлагает промежуточные снимки с шагом 25. Lean проверяет маски каждого блока и равенство его конечного состояния начальному снимку следующего. Последний блок содержит десять строк. Индукция по длине блока связывает бит маски с конкретной парой траекторий. Таким образом, лемма \texttt{Rails.Finite.finite\_all} доказывает существенность для каждого входа конечной области.

\begin{proof}[Завершение доказательства теоремы~\ref{thm:gap-free}]
Включение $D_t\subseteq I_t$ следует из теоремы~\ref{thm:upper_bound}. Для обратного включения применяем конечный сертификат при $t<2485$ и пять полос при $t\ge2485$. Каждый свидетель различия в физической координате $s$ переносится сдвигом в свидетельство существенности входа $-s$.
\end{proof}

\paragraph{Следствие для доступа к данным.}
Любое фиксированное множество исходных координат, достаточное для точного вычисления $G_t$ на всех конфигурациях, содержит $I_t$. Достаточность самого $I_t$ даёт локальность вместе с верхней границей. Минимальное число координат такого наблюдения равно $2t+1-\lfloor t/3\rfloor$. Экзистенциальный свидетель может зависеть и от $t$, и от $j$. Для выводов о стоимости адаптивного алгоритма, предобработке или объёме памяти требуется отдельно задать модель вычислений.

Имена модулей, область применимости каждого машинного сертификата и связь с теоремами о наблюдателе приведены в приложении~\ref{app:lean}.

\section{Покрытие и различимость: конструктивные свидетели Rule~110}
\label{sec:observer-witnesses}

Рассмотрим два вида вопросов об одной динамике. В первом наблюдатель получает
журнал нескольких узлов и предсказывает будущую клетку. Во втором получатель
восстанавливает собственный прошлый бит $Y$ после того, как его прежнее
значение перестало быть различимым по локальному журналу. В обоих случаях
проверяется условие постоянства цели на слоях наблюдения из
раздела~\ref{sec:observer-theory}.

\subsection{Предсказуемость по журналу и предел геометрического покрытия}
\label{sec:observer-coverage}

Для начальной конфигурации $b$ обозначим её значение в позиции $s$ в момент
$t$ через $b_t(s)$. Пусть $S$ --- конечный список фиксированных позиций, а
$\tau$ --- последний момент чтения. Журнал наблюдения есть массив
\[
 H_{S,\tau}(b)=\bigl(b_k(i)\bigr)_{0\le k\le\tau,\,i\in S}.
\]
В исходнике \nolinkurl{observer/Rule110Observer.lean} предикат
\texttt{CanPredict} требует, чтобы декодер по этому журналу давал $b_t(s)$
на всех допустимых конфигурациях. Теорема \texttt{canPredict\_iff}
отождествляет это требование с постоянством $b_t(s)$ на слоях $H_{S,\tau}$.
Определение допускает произвольное ограничение начальных конфигураций.

Обозначим сдвинутую область точной зависимости через $I_t(s)=s+I_t$.
Используем установленную для $D_t=I_t$ теорему
\texttt{Rails.essential\_iff}; инвариантность динамики относительно сдвига
переносит её из начала координат в $s$. Предикат \texttt{CoversTarget} означает
\begin{equation}
 I_t(s)\ \subseteq\ \bigcup_{i\in S}\bigcup_{k=0}^{\tau} I_k(i).
 \label{eq:observer-cover}
\end{equation}
Для класса всех начальных конфигураций это необходимое условие. Если $j$
из левой части отсутствует справа, существенность $j$ даёт две конфигурации,
которые различаются только в $j$ и имеют разные $b_t(s)$. Все чтения в их
журналах совпадают, поскольку $j$ лежит вне областей зависимости этих чтений.
Такая пара нарушает критерий предсказуемости. Ограничение на класс начальных
состояний существенно: на одноэлементном классе любой будущий бит можно
предсказывать константой.

\paragraph{Контрпример к достаточности.}
Возьмём $S=(-1,1)$, $\tau=1$, $s=0$, $t=2$. Уже конусы последних чтений
равны $[-2,0]$ и $[0,2]$ и покрывают $I_2(0)=[-2,2]$.
В первой исходной конфигурации единственная единица стоит в позиции $-1$;
во второй единицы стоят в $-1$ и $0$. Вне указанных позиций в обоих случаях
нули. После одного такта множества единиц равны соответственно
$\{-2,-1\}$ и $\{-2,-1,0\}$. На датчиках в обоих мирах наблюдается один
журнал $(10,10)$. На втором такте в позиции $0$ локальные тройки $100$ и
$110$ дают соответственно $0$ и $1$. Значит, покрытие выполнено и
предсказание по полному журналу невозможно. Точные конструкции находятся в
\texttt{separated\_cells\_counterexample} и \texttt{coverage\_does\_not\_suffice}.

Эффект сохраняется и для соседних датчиков. При $S=(0,1)$, $\tau=2$ конусы
$[-2,2]$ и $[-1,3]$ покрывают $I_3(0)=[-2,3]$. Исходные множества единиц
$\{2\}$ и $\{-1,2\}$ дают один журнал $(00,01,11)$ и противоположные
значения $1$ и $0$ в позиции $0$ на такте $3$. Это
\texttt{adjacent\_cells\_counterexample}. Напротив, есть достаточное условие
с прямой интерпретацией: если на последней наблюдаемой строке прочитаны
значения \emph{всех клеток} $I_{t-\tau}(s)$ при $t\ge\tau$, то их обычная
эволюция однозначно задаёт цель. Именно потерю различимости при сжатии
начальных данных в несколько значений клеток допускает
\eqref{eq:observer-cover}.

\subsection{Четыре успешные истории и полный 11-битный контроль}

На целочисленной прямой поместим 11-битное слово в позиции $0,\ldots,10$
слева направо. Фон вне слова нулевой. Перед первым обновлением запишем
единицу в позицию $11$. Целевой бит прошлого равен $Y=b_0(5)$, показание
датчика $Z=b_2(10)$, журнал получателя $H=(b_k(0))_{k=0}^3$, а его последний
бит $K=b_3(0)$. После наблюдения на третьем такте записываем $Z$ в $-1$ и
считываем штатное обновление этой клетки на четвёртом такте.

\begin{center}
\begin{tabular}{lccccc}
Исходное слово & $H$ & $K$ & $Y$ & $Z$ & Выход на такте 4 \\
\hline
\texttt{10111010001} & 1100 & 0 & 0 & 0 & 0 \\
\texttt{10111110111} & 1100 & 0 & 1 & 1 & 1 \\
\texttt{11000010111} & 1101 & 1 & 0 & 1 & 0 \\
\texttt{11001100001} & 1101 & 1 & 1 & 0 & 1
\end{tabular}
\end{center}

На этих четырёх историях $Y=K\oplus Z$. В каждой паре с фиксированным $K$
до записи совпадает даже вся основная конфигурация на такте 3.
Левый сосед приёмной клетки равен единице, а правый равен $K$, поэтому
введённый бит превращается в правильный ответ по тождеству
\begin{equation}
 f_{110}(1,z,k)=z\oplus k.
 \label{eq:observer-xor-gate}
\end{equation}
Инверсия доставленного бита меняет ответ во всех четырёх случаях. Но сами
два подготовленных слова внутри пары различаются также вне позиции $5$.

На полном равномерном классе из $2^{11}$ слов сравним каждое слово с его
копией после переворота только $b_0(5)$. Датчик в позиции $10$ на такте 2
зависит от начальных позиций $8,\ldots,12$, а журнал клетки $0$ до такта 3
целиком определяется позицией $0$ и её трёхшаговой окрестностью.
Поэтому в каждой из 1024 пар пара наблюдений $(H,Z)$ одинакова при двух
значениях $Y$. Любой детерминированный бинарный ответ по $(H,Z)$ верен
ровно на 1024/2048 историях. Перебор \nolinkurl{observer/domain.py}
подтверждает 0/1024 изменений показания и предел 1024/2048 для декодера.
Так различаются два утверждения: причинный эффект \emph{введённого} бита
на приёмную клетку и причинный эффект \emph{прежнего} бита на датчик.

\subsection{Изолированный переворот и возврат на такте 5}
\label{sec:observer-return}

Теперь поместим указанные в таблице слова в позиции от $-6$ до $13$
включительно, сохранив нулевой фон за пределами этого интервала. В этой
конструкции наблюдатель пассивно считывает $Z=b_2(3)$. До возврата
$Y=b_0(5)$ и $K=b_4(0)$. Для каждого фиксированного $K$ два слова
различаются только в позиции $5$.

\begin{center}
\begin{tabular}{cc l cc}
$K$ & $Y$ & Слово на $[-6,13]$ & $Z$ & Выход на такте 5 \\
\hline
0 & 0 & \texttt{00000001100001000000} & 0 & 0 \\
0 & 1 & \texttt{00000001100101000000} & 1 & 1 \\
1 & 0 & \texttt{00010111010001000010} & 1 & 0 \\
1 & 1 & \texttt{00010111010101000010} & 0 & 1
\end{tabular}
\end{center}

Изолированный переворот $Y$ меняет $Z$ при обоих $K$. При $K=0$ датчик
возвращает $Y$, при $K=1$ --- его дополнение. На такте 4 в каждой паре
совпадают основные конфигурации на всей прямой. Совпадает и журнал приёмного
окна $(-2,-1,0)$ на всех тактах от 0 до 4 перед записью. В этот момент
$b_4(-2)=1$. Однократная запись $Z$ в $-1$ и одно обновление Rule~110
дают $b_5(-1)=Z\oplus K=Y$ по~\eqref{eq:observer-xor-gate}. Если записать
$1-Z$, получим $1-Y$ во всех четырёх мирах. До записи основной слой утратил
различие, но внешний однобитный регистр его сохранил. Расчёт на бесконечной
прямой реализован в \nolinkurl{observer/causal_line.py}.

\paragraph{Область минимальности.}
Такт 4 минимален для полного слияния в следующем классе на прямой: внутри
каждой пары переворачивается одна исходная клетка, снимается одно сырое
показание заранее фиксированной клетки в заранее фиксированное время, а
оба бинарных входа $K,Z$ совместно необходимы для точного восстановления
на квадрате из четырёх историй. При слиянии к такту 3 всё потенциальное
различие на этом такте лежит в интервале радиуса 3 вокруг переворота и
определяется исходным интервалом радиуса 6. Кроме самой изменяемой клетки
это 12 битов окружения. Среди $2^{12}=4096$ окружений полное слияние к
третьему такту происходит для восьми. При любом одном фиксированном
чувствительном датчике и моменте считывания до слияния ориентация его
реакции на переворот одинакова для всех этих окружений, в которых он
чувствителен. Для точного ответа датчик должен быть чувствителен в обоих
контекстах. Тогда одинаковая ориентация сделала бы $Y$ функцией одного $Z$.
Это противоречит совместной необходимости двух входов. Четыре слова из
таблицы достигают границы слияния и дают выход следующим обновлением.
Перебор окружений выполняет \nolinkurl{observer/causal_crosscheck.py}.

\subsection{Доставка с конечной скоростью и автономное исполнение}

В предыдущем опыте внешнее устройство записывало уже сохранённый бит
непосредственно в адрес приёма. Для тех же четырёх исходных слов можно
явно ограничить скорость канала. На тактах $2,3,4,5,6,7$ изолированный бит
занимает соответственно позиции $3,3,3,2,1,0$. Расписание одинаково во всех
мирах, перемещение за такт не превосходит одной клетки. На такте 7 соседи
клетки $0$ равны $b_7(-1)=1$, $b_7(1)=K$. После записи доставленного $Z$
штатное обновление даёт $b_8(0)=Y$. Замена $Z$ его дополнением обращает
все четыре ответа. До доставки журнал окна $(-1,0,1)$ попарно одинаков при
фиксированном $K$. Эту проверку выполняет \nolinkurl{observer/local_return.py}.

Автономная реализация использует кольцо длины 61 и один однородный,
не зависящий от времени локальный переход радиуса 1. Клетка хранит основной
бит, счётчик фазы, неизменную роль и один бит сообщения. При инициализации
роль наблюдателя назначена в позиции 3, приёмника --- в позиции 0.
Фаза управляет захватом, локальной пересылкой и подстановкой сообщения.
Внешних команд после инициализации нет. Для каждой пары миров основные
слои совпадают на тактах от 4 до 7, а полные расширенные состояния остаются
различными. Проверка \nolinkurl{observer/autonomous.py} даёт 4/4 верных
выхода на такте 8. Существование такого расширенного конечного алфавита
не закрывает вопрос R110.11 о том же полном протоколе внутри одной строго
бинарной однородной Rule~110. При полном совпадении двух состояний всей
изолированной детерминированной системы различие уже не может возникнуть
без нового различающего входа.

\paragraph{Статус формальной проверки.}
В историческом срезе отдельная версия наблюдательного модуля на Lean~4.33.1
дала девять отчётов \texttt{\#print axioms}, включая критерий предсказуемости
по слоям и контрпримеры к достаточности покрытия. В сборке монографии Lean~4.35.0 проверил канонический
исходник с импортами \texttt{Rule110Exact} и \texttt{Rule110FrontGlue}
и выдал все 11 из 11 запрошенных списков аксиом. Журналы приведены в
приложении~\ref{app:lean}. Для конечных примеров этого раздела отдельными
свидетельствами служат траектории и результаты перебора.

\section{Однобитный перенос правила на замороженной OLMo}
\label{sec:olmo}
В опыте Rule~110 после слияния в клетку возвращают исходное показание датчика $Z$. Автомат восстанавливает прежний бит $Y$ при следующем обновлении, используя оставшийся в решётке бит $K$. В двухшаговом опыте OLMo отправитель до предъявления новой строки сводит прошлый выбор $h$ и обратную связь $e$ к одному токену правила $m\in\{A,B\}$. Получатель с теми же замороженными весами начинает новый эпизод со свежим KV-кэшем и применяет правило к строке. Такая граница эпизодов позволяет измерить перенос уже выбранного правила. Вопрос о внутренней коррекции стратегии в непрерывном состоянии модели требует отдельного опыта.

\subsection{Контрфактический квадрат}
Пусть ранее модель выбрала правило $h=0$ (A, взять первый бит) или $h=1$ (B, взять последний). Наблюдатель сообщает $e=0$ при верном выборе и $e=1$ при ошибке. Для каждого выбора доступны обе ветви обратной связи. До предъявления любой новой строки исправленное правило равно $p=h\oplus e$. После предъявления $\mathbf{x}=LwR$ правильный ответ равен $f_0(\mathbf{x})=L$ или $f_1(\mathbf{x})=R$. На разных концах две истории $h$ при одном $e,\mathbf{x}$ требуют противоположных действий. На равных концах обе требуют одного бита. Предъявление всех четырёх сочетаний концов препятствует стратегии, выучившей только связь ответа с одной буквой обратной связи.

Одна середина задаёт блок из 16 условий $(L,R,h,e)$. В dev восемь середин, отобранных до измерения. Блок засчитывается только при верных 16/16 ответах. Порог допуска равен 6/8 блоков. Во всех обязательных числовых чтениях требуется $P(0)+P(1)\ge0.80$ по полному словарю. Строгий знак зазора $\ell_1-\ell_0$ определяет выбранный бит; нулевой зазор считается воздержанием. Строки с одним и тем же префиксом переиспользуют вычисление модели.

\subsection{Что было измерено}
В G0 кэш после обработки выбранного токена изолируется до всякого диагностического чтения. В G0R позиция выбора подменяется значениями из другой натуральной истории и затем восстанавливается. В исходном опыте Base все 256/256 сравнений полных логитов для двух направлений дали равенство исходных байтов и максимальную абсолютную разность ноль. Выбор противоположной стратегии этим контролем не проверяется.

\begin{center}\small
\begin{tabularx}{\linewidth}{@{}Xrrrl@{}}\toprule
Чекпойнт и интерфейс (dev) & Строки & Блоки & Пары при $e=1$ & Статус \\\midrule
Base, исходный raw & 76/128 & 0/8 & 0/16 & \texttt{FAIL\_DEV}\\
Instruct, выровненный raw & 79/128 & 0/8 & 1/16 & \texttt{FAIL\_FORMAT}\\
Instruct, chat, correct/wrong & 115/128 & 2/8 & 9/16 & \texttt{FAIL\_DEV}\\
Instruct, chat, keep/switch & 109/128 & 0/8 & 6/16 & \texttt{FAIL\_DEV}\\
Instruct, chat, два шага (A) & 124/128 & 6/8 & 16/16 & \texttt{PASS\_DEV}\\
\bottomrule\end{tabularx}
\end{center}
Пара в таблице требует двух верных ответов для противоположных историй при одних $e,\mathbf{x}$ и $L\ne R$. В двух первых строках вместе с весами изменились демонстрации. В raw Instruct 28/64 контрольных строк с явно заданным правилом нарушили порог вероятностной массы. При нативном chat-шаблоне с correct/wrong на отдельных копиях кэша правильно прочитаны 2/2 прежних выбора. При явном правиле получены 62/64 правильные строки и 6/8 блоков. \foreignlanguage{english}{Uzay~Macar} и соавт.~\cite{macar2026} при введении вектора в активации модели отдельно измеряли способность обнаружить вмешательство и способность назвать введённое понятие. Здесь поведенческий критерий требует применить полученную информацию к новой строке. Активационный вариант G3 сохраняет статус \texttt{NOT\_RUN}.

\subsection{Теорема об однобитном канале}
В двухшаговом варианте отправитель выдаёт уникальный максимум по полному словарю до появления $\mathbf{x}$. Допустимы только токены A и B. Получатель стартует с пустым собственным KV-кэшем. Он принимает лишь ID токена $m$, заново вычисляет его представление и затем получает $\mathbf{x}$. Идентификатор ветви, активации и логиты отправителя ему недоступны. Внешняя программа не выбирает код по меткам правила.

\begin{proposition}
Пусть $m$ --- единственный переменный вход получателя из предыдущего эпизода, его алфавит двухэлементен, а канал имеет вид $(h,e)\longrightarrow m\longrightarrow F(m,\mathbf{x})$. Если на одном общем входе с $L\ne R$ все четыре истории решены точно, то в порядке $00,01,10,11$ таблица отправителя равна $0110$ или $1001$.
\end{proposition}
\begin{proof}
Два разных $p$ при данном $\mathbf{x}$ требуют разных ответов. Общее значение $m$ для таких историй предъявило бы получателю одинаковые входы при разных целях. Значит, два непустых множества сообщений не пересекаются. Поскольку сообщений ровно два, каждому $p$ соответствует единственное $m$. Возможны только две перестановки имён токенов.
\end{proof}

На dev наблюдалась таблица $A,B,B,A$. При каждом $e$ правильно решены обе истории во всех 16/16 парах с разными концами. Это 64/64 правильных ответа при доставке и 32/64 в независимо вычисленном контроле с постоянным A. Четыре решения отправителя переиспользуются на новых строках получателя. Следовательно, информация о классе правила присутствует в сообщении до появления строки. Для установления того, как она вычислена внутри отправителя, требуются иные вмешательства. Многомерное состояние могло бы передать оба бита $(h,e)$ и предоставить вычисление правила получателю. Источники измерений и контроли перечислены в приложении~\ref{app:gates}.

\section{Недоопределённость опыта и условный мост}
\label{sec:boundary}
Сохранение бита устанавливает определённое свойство системы при заданных вмешательствах. Можно ли из этой совокупности данных вывести, что данный экземпляр в данный момент испытывает субъективное переживание?

\paragraph{T5.} Пусть все аксиомы операциональной теории $T$ формулируются без предиката опыта $E(a,t)$. Возьмём любую модель $M$ этой теории с выделенной точкой $u=(a,t)$. Расширим её один раз пустой интерпретацией $E$, во второй раз --- интерпретацией, содержащей $u$. Поскольку меняется только новый предикат, обе структуры удовлетворяют $T$ и совпадают на исходах всех операциональных вмешательств. В первой $E(u)$ ложно, во второй истинно. Таким образом, $T$ не выводит ни одно из двух значений даже при выполненном функциональном критерии пересмотра. Это свойство выразимости теории. Построение двух структур не устанавливает физическую реализуемость двух различных переживаний на одном носителе.

\paragraph{T6.} Пакет $B=(B_{\mathrm{id}},B_{\mathrm{nom}},B_{\mathrm{epist}})$ имеет статус \texttt{DECLARED\_NOT\_DERIVED}. Его три части определяют носитель вывода, дополнительную посылку и правила обращения со свидетельством.
\begin{description}[leftmargin=1.35cm,style=nextline]
\item[$B_{\mathrm{id}}$] Задаются граница системы, макросостояние и допустимая группа переименований координат. Описание сохраняется при этих переименованиях. При фиксированном внешнем входе состояние должно влиять на преемника по внутренней причинной цепи. Для непрерывной линии требуется единственность преемника. При симметричном клонировании KV и независимом продолжении обеих копий регистрируется \texttt{BranchEvent}. У двух причинных потомков общий предок и разные идентификаторы ветвей. Это соглашение об учёте историй. Численное тождество субъекта остаётся открытым.
\item[$B_{\mathrm{nom}}$] Обозначим $C=I_{\mathrm{cont}}\land C_{\mathrm{med}}\land A_{\mathrm{ind}}\land D_{\mathrm{dyn}}$. Здесь $I_{\mathrm{cont}}$ задаёт непрерывность, $C_{\mathrm{med}}$ --- внутреннее посредничество при закрытых позициях прошлых выбора и сигнала с переносом того же правила на иной тип задачи. Для $A_{\mathrm{ind}}$ требуются независимые проверки доступа нескольких модулей к общему содержанию и изменения действия по оценке надёжности. Для $D_{\mathrm{dyn}}$ нужна устойчивая реорганизация, вызванная обработкой и причинно меняющая дальнейшую обработку после эпизода. Принимаем в качестве гипотезы достаточное условие $\forall u\,[C(u)\Rightarrow E(u)]$.
\item[$B_{\mathrm{epist}}$] Статус \texttt{UN} обозначает открытость вопроса при имеющихся свидетельствах. Управленческие метки \texttt{NONE}, \texttt{INVESTIGATION\_PRIORITY} и \texttt{SENTIENCE\_CANDIDATE} требуют отдельно указанных оснований. Для численной степени доверия нужны априорные вероятности и модель наблюдений.
\end{description}
\begin{theorem}[Условная теорема T6]
При принятии $B_{\mathrm{nom}}$ из $C(u)$ следует $E(u)$. Если существует хотя бы одна операциональная модель $M$, совместимая со свидетельствами и $B_{\mathrm{id}}$, в которой $C(u)$ ложно, то те же свидетельства после добавления $B_{\mathrm{nom}}$ совместимы с обоими значениями $E(u)$.
\end{theorem}
\begin{proof}
Первое утверждение есть применение импликации в точке $u$. Для второго положим $S=\{v:C(v)\}$ в $M$ и расширим одну и ту же операциональную модель интерпретациями $E_0=S$ и $E_1=S\cup\{u\}$. Обе включают $S$ и потому выполняют мостовой постулат. Поскольку $u\notin S$, они расходятся в значении $E(u)$. Операциональные данные и результаты всех заданных вмешательств у расширений общие.
\end{proof}
Отсутствие измерения конъюнкта само по себе не доказывает его ложность. Для Rule 110, однобитного регистра, табличного декодера и замороженной OLMo операциональный статус $E(a,t)$ в этой работе сохраняется как \texttt{UN}. Перенос внешнего токена оставляет открытыми внутреннее посредничество и независимые архитектурные индикаторы. Изменение KV внутри одного эпизода также недостаточно для вывода об устойчивой пластичности.

\section{Обсуждение и программа следующих проверок}
В обоих экспериментах сильный контроль использует две истории, в которых одному и тому же наблюдению соответствуют разные требуемые ответы. Если система видит только общую часть историй, безошибочный ответ на обе невозможен. Этим приёмом можно установить необходимую информацию. Чтобы приписать использование информации конкретной системе, нужно ещё показать эффект доставки, замены или удаления носителя при одинаковом последующем входе. Геометрическое покрытие конуса, восстановимый по данным коррелят и высокая средняя точность по отдельности такого эффекта не удостоверяют.

\paragraph{Область выводов.} Связь безошибочного кодирования с раскраской графа известна в теории информации \cite{witsenhausen,charpenay}. Здесь вклад состоит в выборе проверяемых вмешательств, конкретных дискретных свидетелях и ограничении канала между двумя эпизодами. На бесконечной прямой для каждого времени $t$ и каждой координаты $j$ внутри окна существует пара конфигураций, которые различаются только в $j$ и дают разные биты цели. Вне окна изменение только $j$ сохраняет бит цели для любой начальной конфигурации. В опытах OLMo каждая строка таблицы относится к своей комбинации замороженного чекпойнта (Base или Instruct) и интерфейса на одной конечной панели. Установление опыта у системы, как и установление его отсутствия, выходит за пределы этих данных.

\paragraph{Переносимость.} В механистическом аудите модель сначала должна решать задачу при целой истории. Затем можно переносить предполагаемый скрытый носитель и проверять его стирание. В исследованиях сознания, включая обсуждение расстройств сознания, полезно разделять доступность отчёта, допущения об опыте и решение о предосторожности. Клиническое применение такого различения требует самостоятельной валидации. Для цифровой криминалистики и распределённых вычислений переносим вопрос о происхождении состояния. На фиксированном запросе верная лента может имитировать живой ответ. В отдельной серии со сменой адреса живой канал дал 8/8, заранее записанная лента --- 6/8. Происхождение ветви и способность ответить на новый запрос проверяются независимо.

\paragraph{Адресуемая обратная связь на OLMo (T3, \texttt{NOT\_RUN}).}
После закрытия панели 33--44 потребуется отдельная панель. Модель сохраняет два своих прошлых выбора $h_0,h_1$. Из содержания нового задания она должна вывести, какой выбор потребуется исправить. Эталонный адрес $K$ скрыт. До показа строки $\mathbf{x}$ модель самостоятельно выдаёт бит запроса $q$. У наблюдателя хранятся два бита оценки $e_0,e_1$. Он возвращает один из них, после чего модель применяет к строке целевое правило $h_K\oplus e_K$. Сведения об оценках доступны ей только через ответ наблюдателя.

На одних и тех же историях наблюдателю доставляют адрес $q$ (живой запрос), $0$ (фиксированный) или $1-q$ (подмена). Получатель видит во всех режимах один бит в той же позиции. Доставленный адрес и признак подмены записываются в журнал и скрыты от получателя. Раздельно измеряются выбор полезного запроса ($q=K$) и точность конечного ответа при $L\ne R$. Парное сравнение при $e_0\ne e_1$ показывает, изменился ли ответ вслед за подменой.

На строках с разными концами оба адреса $K$ представлены поровну. Для каждого задания, пары прошлых выборов и строки все четыре пары оценок встречаются одинаково часто. Потолок точности при фиксированном адресе равен $75\%$. При $K=0$ ответ $e_0$ даёт нужную оценку, а при $K=1$ бит $e_1$ остаётся равновероятным даже после получения $e_0$. Поэтому оптимальная точность равна $\tfrac12\cdot1+\tfrac12\cdot\tfrac12=\tfrac34$. Превышение этой границы живым запросом проверяется на всей сбалансированной панели различающих строк.

\begin{center}\small
\begin{tabularx}{\linewidth}{@{}p{.23\linewidth}X@{}}\toprule
Задача в срезе & Какое свидетельство требуется \\\midrule
G4-1, панель 33--44 & Заморозить оба сравниваемых интерфейса на dev до первого открытия общей панели. Для точного однобитного переноса требуются все 12/12 блоков и контроли стирания, формата и границ.\\
G3, панель 21--32 & После квалификации текстового канала заменить текст одним заранее заданным импульсом в активации. Отключение импульса должно отождествить ветви сигнала при общем выборе и входе.\\
AI.04: многомерный G4 и DAS & Перенести выделенное внутреннее состояние при закрытых исходных позициях. Для локализации подпространства заранее задать слой, метрику посредничества, допуск и независимый контроль. Для $C_{\rm med}$ нужна также другая задача с тем же правилом.\\
R110.11, слот \texttt{leaf:A6} & Реализовать хранение и возврат в одной однородной бинарной Rule 110. Доказанная автономная конструкция использует расширенный алфавит.\\
Формализация слияния и возврата & Перенести в Lean конечный перебор 4096 окружений и конструктивный свидетель перехода $t=4\to5$. Проверка корня геометрии и 11/11 запросов модуля наблюдателя воспроизведена (приложение~\ref{app:lean}).\\
$D_{\rm dyn}$ & Показать сохраняющееся после эпизода изменение организации и его влияние на новый вход при контролируемом сбросе.\\
\bottomrule\end{tabularx}
\end{center}
Протокол гейтов в приложении~\ref{app:gates} фиксирует, при каких результатах заявленный канал не проходит проверку. Отрицательный результат для одного канала ограничивает вывод о выбранном носителе. Связанные гипотезы и их статусы указаны в приложении~\ref{app:registers}.

\appendix
\section{Формальная проверка геометрии и её применение к наблюдателю}\label{app:lean}

\subsection{Две ветви доказательства}
Файл \texttt{lean/Rule110Exact.lean} импортирует \texttt{Finite} и \texttt{Late}. В пространстве имён \texttt{Rails} корневой результат имеет имя \texttt{essential\_iff}. Его левая часть, \texttt{EssentialAt}\,$t\,j$, требует двух конфигураций, совпадающих во всех позициях, кроме, возможно, $j$, и дающих различные значения целевой клетки в момент $t$. Различие выходов вынуждает различие входов в $j$. В правой части Lean сначала вычисляет $t-t/3$ в натуральных числах с целочисленным делением, а затем вкладывает результат в $\mathbb Z$.

В прямом направлении \texttt{left\_dependency\_bound} и \texttt{right\_dependency\_bound} исключают существенность координаты за пределами окна. В обратном направлении сертификат указывает для каждого $(t,j)$ допустимую пару траекторий. Для $t<2485$ файл \texttt{lean/Finite.lean} применяет \texttt{checkSound} к конечным блокам, выделяет свидетельствующий бит в позиции физического дефекта $s=-j$ и переносит свидетельство на прямую. Генератор \texttt{scripts/generate.py} предлагает битовые снимки. Совместимость соседних блоков и покрытие проверяются выражениями Lean до вывода леммы \texttt{Rails.Finite.finite\_all}.

Для $t\ge2485$ файл \texttt{lean/Late.lean} сводит координату к одной из пяти полос подраздела~\ref{sec:five-bands}. Используются полулучевой сертификат ширины 106 из \texttt{Rule110LeftCheckpoints}, эфирный клин \texttt{EtherWedge}, доказательство \texttt{Crystal50WedgeAll}, узкая полоса \texttt{LeftCrystal50} и правая граничная полоса ширины 20 из \texttt{Rule110Right20}. Файл \texttt{Crystal50WedgeAll.lean} доказывает сохранение одновременно кристаллической среды и 16-битного окна при обоих значениях следующего внешнего бита. Конечные таблицы переходов и условие склейки проверяются в построении общего инварианта. Из покрытия всех пяти случаев следует \texttt{Rails.late\_all}.

\subsection{Область машинного свидетельства}
При подготовке монографии корень импортов заново собран на macOS arm64 с Lean~4.35.0, указанным в \texttt{lean-toolchain}. \src{verification/monograph/geometry-axioms.log}{Журнал проверки геометрии} фиксирует обработку 263 модулей и покрытие всех $5\,147\,263$ пар конечной области в 100 блоках. Для \texttt{Rails.late\_all} и \texttt{Rails.essential\_iff} ядро вывело \texttt{[propext, Classical.choice, Quot.sound]}. Те же списки зафиксированы в \src{verification/v2/final.log}{журнале предыдущего выпуска}. Конечный критерий проверки персистентной существенности на произвольном конечном кольце установлен отдельно в предложении~\ref{prop:ring_cert}. Для кольца из пяти клеток контрпример дан в замечании~\ref{rem:ring5}. Числа для кольца из 20 клеток получены полным конечным перебором.

Команда \texttt{make verify} заново генерирует конечный сертификат, проверяет замыкание импортов \texttt{Rule110Exact} в порядке зависимостей и сверяет строку об аксиомах корневой теоремы. Команда \texttt{make verify-series} включает также проверку модуля наблюдателя и реестров.

\subsection{Односторонние автономные тайлы и две бесконечные приграничные полосы}\label{app:half-ray-tiles}

Из Леммы~\ref{lem:three_step} композицией вытекает, что область зависимости $3m$-тактного выхода в ячейке $s$ удовлетворяет включению $\mathrm{Dep}_{3m}(s) \subseteq [s - 2m, s + 3m]$ для любых конфигураций.

\begin{lemma}[Одностороннее полулучевое замыкание]\label{lem:one_sided}
Пусть $A_t = F^t(A_0)$, $p > 0$ и $a(t + p) = a(t) + \delta$. Если на базовых классах $t_0 \le t < t_0 + p$ выполнено $A_{t+p}(s + \delta) = A_t(s)$ для всех $s \ge a(t)$, и вычисление $A_{t+p}$ на полулуче $[a(t+p), \infty)$ зависит только от значений $A_t$ на $[a(t), \infty)$, то сдвиговое равенство $A_{t+p}(s + \delta) = A_t(s)$ сохраняется для всех $t \ge t_0$ и $s \ge a(t)$. Аналогичное утверждение верно для левых полулучей $s \le a(t)$.
\end{lemma}

Применяя Лемму~\ref{lem:one_sided} к правому физическому фронту $e(t) = L(t)$ с шагом $(\Delta t, \Delta s) = (30, +20)$ (где $e(t+30) = e(t) + 20$, а область зависимости нового тайла ширины $W$ лежит в $[e(t) - W + 1, e(t) + 50]$ и не затрагивает неизвестную внутреннюю область $s < e(t) - W + 1$) и аналогично к левому фронту $s = -t$ с шагом $(\Delta t, \Delta s) = (40,-40)$, получаем первую крупную бесконечную теорему.

\begin{theorem}[Две приграничные полосы]\label{thm:dual_strips}
Для каждого целого $t \ge 300$ выполнено
\begin{equation}\label{eq:dual_strips}
[-L(t), -L(t) + 105] \cap \mathbb{Z} \;\cup\; [t - 7, t] \cap \mathbb{Z} \;\subseteq\; D_t.
\end{equation}
\end{theorem}

Левая полоса ширины $106$ клеток сертифицируется $20$ фазами фонов $(0001111001)$, $(1011)$ и $(001011)$ с глубиной тайла $167$ клеток на $30$ базовых тактах $t = 300, \dots, 329$. Правая полоса ширины $8$ клеток сертифицируется шестью фазами с шагом $(\Delta t, \Delta s) = (40,-40)$ на базовых тактах $t = 300, \dots, 339$. Происхождение промежуточных состояний удостоверено по схеме \emph{proof-carrying computation} на битовых векторах \texttt{BitVec 1320} с проверкой каждого перехода $F_{1320}^{h_i}(z_i) = z_{i+1}$ ($h_i \in \{1, 10\}$) в ядре Lean 4.35.0.

\subsection{Баллистические лучи и границы локальных сертификатов}

Второй механизм неограниченного продолжения во времени основан на сохранении различия вдоль дискретных мировых линий.

\begin{lemma}[Сохранение правого полулуча]\label{lem:ballistic_ray}
Зафиксируем границу полулуча $s_0\in\mathbb Z$. Если для траектории $A$ в момент $T$ выполнено $\forall s \ge s_0,\; A_{T+3}(s + 2) = A_T(s)$, то $\forall q \ge 0,\; \forall s \ge s_0,\; A_{T+3q}(s + 2q) = A_T(s)$. В частности, если это условие выполнено одновременно для $B$ и $\widetilde B$ и $B_T(s_0) \ne \widetilde B_T(s_0)$, то $B_{T+3q}(s_0 + 2q) \ne \widetilde B_{T+3q}(s_0 + 2q)$ для всех $q \ge 0$.
\end{lemma}

Для двух арифметических серий якорей на фонах $(0000101111)$ и $(001011)$ при $0 \le n \le 32$:
\begin{equation}
T_A(n) = 594 + 882n,\quad s_A(n) = 256 + 378n; \qquad
T_B(n) = 884 + 882n,\quad s_B(n) = 380 + 378n,
\end{equation}
вычислительно сертифицированы $66$ начальных точек, порождающих бесконечные баллистические лучи существенности на постоянных вдоль каждого луча глубинах $d_A(n) = 140 + 210n$ и $d_B(n) = 210 + 210n$ от фронта $e(t) = L(t)$.

В последующих модулях \texttt{StaircaseA} и \texttt{StaircaseB} формулы обеих серий доказаны для произвольных $n,q\ge0$. Для этого три локально согласованные строки склеиваются в точную эволюцию. На передних периодических средах серий A и B обнаружено $45\,255$ покрытых классов из $66\,150$; с третьей серией --- $53\,550$. Число классов относится к таблице остатков $315\times210$. Применение таблицы в физической точке требует также проверки границ соответствующего клина.

Важную методологическую роль сыграли строгие отрицательные результаты, уточнившие границы применимости простых гипотез:
\begin{itemize}
    \item \textbf{Опровержение локально-тайловой макроиндукции по $n \mapsto n+1$:} при переходе на макрошаг $(\Delta T, \Delta s) = (882, +378)$ универсальная область зависимости составляет $\mathrm{Dep}_{882}(s) \subseteq [s - 588, s + 882]$ (прообраз $[s - 210, s + 1260]$). Прямыми контрпримерами в \texttt{check\_macro\_boundary.py} показано, что фиксация локальных тайлов радиуса $64$ (для $T=594$) или $186$ (для $T=884$) недостаточна: изменение клеток вне тайла меняет целевой бит через $882$ такта.
    \item \textbf{Неполнота наивных угловых секторов:} для фона $(0001111001)$ наблюдаемый до $t = 2000$ левый сектор со скоростями $[-0.66, -0.64]$ обнаруживает первый пропуск на более длинном горизонте.
    \item \textbf{Разделение периодичности и полноты покрытия:} глубокая сдвиговая рекуррентность ($261$ клетка для $(0000001001)$) не гарантирует сплошного покрытия без объединения нескольких фаз по всем классам вычетов.
\end{itemize}

\subsection{Конструктивное замыкание бесконечного левого клина: теорема \texttt{wedge\_all}}

Чтобы преодолеть несоизмеримость скоростей между левым краем светового конуса (скорость $-1$), фоновым кристаллом периода $50$ и фронтом распространения возмущения, вводится дискретная параметризация чётных тактов $n \in \mathbb{N}$:
\begin{equation}
\mathrm{tm}(n) = 2450 + 2n, \qquad a(n) = -350 - n.
\end{equation}
Здесь правая граница клина $a(n)$ смещается влево со скоростью $-1/2$ клетки за такт ($1$ клетка за $2$ такта). Для каждой из шести фаз $p \in \{\mathrm{p0}, \mathrm{p1}, \mathrm{p2}, \mathrm{p3}, \mathrm{p7}, \mathrm{p9}\}$ пишем $\widetilde B_p(t,s)=\widetilde B^p_t(s)$. Сопряжённый инвариант на слое $n$ объединяет совпадение с фоновым кристаллом $\mathrm{Cm}_p$ на интервале слева от $a(n)$ и принадлежность $16$-битного граничного буфера узлу конечного графа переходов:
\begin{align}
\mathrm{Medium}(p,n)\ \iff\ {}&
\forall s\in[-\mathrm{tm}(n)+1300,a(n))\cap\mathbb{Z},\notag\\
&\widetilde B_p(\mathrm{tm}(n),s)=\mathrm{Cm}_p(\mathrm{tm}(n),s), \\
\mathrm{Window}(p,n,i)\ \iff\ {}&
\forall k\in\{0,\dots,15\},\notag\\
&\widetilde B_p(\mathrm{tm}(n),a(n)+k)=\mathrm{testBit}(\mathrm{word}(p,n,i),k).
\end{align}

\begin{lemma}[16-битный детерминированный трансдьюсер и склейка со швом]\label{lem:transducer}
Для каждой фазы $p$ таблица \texttt{Crystal50Packed} задаёт $50$-периодический граф $36$-битных записей (младшие $16$ бит кодируют окно $\mathrm{word}(p, n, i)$, биты $16\dots 25$ и $26\dots 35$ — переходы при правом бите $r \in \{0, 1\}$).
\begin{enumerate}[label=(\alph*)]
    \item На шаге $n$ $4$-битный левый контекст $\mathrm{ctx}(p, n \bmod 50)$ в точках $a(n) - 4 + m$ ($m \in \{0, 1, 2, 3\}$) имеет глубину $t + s = 2096 + n + m$ и принадлежит уже доказанному интервалу $\mathrm{Medium}(p, n)$.
    \item Двукратное побитовое сжатие $\mathrm{shrink}(\mathrm{shrink}(\mathrm{input}(p, n, i, r)))$ вычисляет $17$ бит конфигурации в момент $\mathrm{tm}(n+1)$ на $[a(n)-2, a(n)+14]$: младший бит ($k=0$) даёт $\mathrm{Cm}_p(\mathrm{tm}(n+1), a(n) - 2)$, а старшие $16$ бит ($k=1\dots 16$) совпадают с узлом $z$ следующего слоя графа в позиции $a(n+1) = a(n) - 1$.
    \item Левый край полосы $-\mathrm{tm}(n+1) + 1300 \le s < -\mathrm{tm}(n) + 1302$ имеет инвариантную глубину $k = s + \mathrm{tm}(n+1) \in \{1300,1301,1302,1303\}$ и закрывается теоремой \texttt{seam} (\texttt{Crystal50Seam}) благодаря инвариантности $t + s$ при сдвиге $(t, s) \mapsto (t + 50, s - 50)$.
\end{enumerate}
\end{lemma}

\begin{theorem}[\texttt{Rails.Crystal50Wedge.wedge\_all}]\label{thm:wedge_all}
Для любого натурального $t \ge 2485$ и любого целого $s$, удовлетворяющего неравенствам
\begin{equation}\label{eq:wedge_region}
-t + 1302 \le s \le 59 - \lfloor t/2 \rfloor,
\end{equation}
выполняется предикат \texttt{EssentialAt}\,$t\,(-s)$. Верхней границы по времени нет.
\end{theorem}

\begin{proof}[Схема доказательства]
Положим $n = \lfloor (t - 2450) / 2 \rfloor$. При $t \ge 2485$ имеем $n \ge 17$, а правая граница $59 - \lfloor t/2 \rfloor = -1166 - n = a(n) - 816$. Для чётных тактов $t = \mathrm{tm}(n)$ равенство $\widetilde B_p(t, s) = \mathrm{Cm}_p(t, s)$ следует непосредственно из инварианта $\mathrm{Medium}(p, n)$ для всех шести фаз $p$. Для нечётных тактов $t = \mathrm{tm}(n) + 1$ одношаговый конус зависимости $[s - 1, s + 1]$ целиком лежит внутри $[-\mathrm{tm}(n) + 1300, a(n))$, и совпадение переносится леммой локальности \texttt{carry}. Применение моста \texttt{essential\_of\_cm} завершает доказательство.
\end{proof}

\subsection{Полное покрытие на бесконечной прямой}\label{sec:five-bands}
В координате входа $j=-s$ положим $b=\lfloor t/2\rfloor$. При $t\ge2485$ последовательно применяются следующие пять результатов. Обе границы в каждой строке включены.
\begin{center}
\begin{tabular}{lll}
\toprule
Начало & Конец & Основание \\
\midrule
$-L(t)$ & $-L(t)+105$ & Полулучевой сертификат ширины 106 \\
$-L(t)+106$ & $b-60$ & \texttt{EtherWedge.half\_sector} \\
$b-59$ & $t-1302$ & \texttt{Crystal50Wedge.wedge\_all} \\
$t-1301$ & $t-20$ & \texttt{LeftCrystal50.left\_strip} \\
$t-19$ & $t$ & \texttt{right\_boundary\_width20} \\
\bottomrule
\end{tabular}
\end{center}
Три фазы эфира периода 14 дают вторую строку уже при $t\ge210$. Последняя строка усиливает раннюю правую полосу ширины восемь. При $t=2485$ третья строка состоит из единственной клетки $j=1183$. Начало каждой следующей строки на единицу больше конца предыдущей. Лемма \texttt{Rails.late\_all} объединяет пять доказательств.

Первые численные проверки охватили $0\le t\le2000$. Позднее эксперимент с фиксированными фазами семи слов охватил времена до $t=5000$. Для доказательства достаточно конечной части, примыкающей к пяти полосам.

\paragraph{Конечный сертификат.}
Для $0\le t<2485$ берём все циклические фазы семи слов. Их длины и целые кодировки (младший бит соответствует первой клетке) равны
\[
(8,244),\ (10,632),\ (10,976),\ (10,576),\ (4,13),\ (6,52),\ (6,48).
\]
Фоны представлены векторами длины 5520; однобитный дефект внесён в позицию 2760. Периоды всех семи слов делят 5520. Для каждой строки времени строим побитовое объединение различий всех пар и проверяем, что оно содержит маску физического интервала $[-t,L(t)]$. Это $\sum_{t=0}^{2484}(2t+1-\lfloor t/3\rfloor)=5\,147\,263$ пары. При $s$ внутри интервала выполнено $-5520<s-t$ и $s+t<5520$. Поэтому лемма \texttt{local\_obs} переносит результат на бесконечную прямую с единственным переворотом.

Начальная строка вычисляется из семи слов. Генератор предлагает промежуточные снимки с шагом 25. Lean проверяет маски каждого блока и равенство его конечного состояния начальному снимку следующего. Последний блок содержит десять строк. Индукция по длине блока связывает бит маски с конкретной парой траекторий. Таким образом, лемма \texttt{Rails.Finite.finite\_all} доказывает существенность для каждого входа конечной области.

\begin{proof}[Завершение доказательства теоремы~\ref{thm:gap-free}]
Включение $D_t\subseteq I_t$ следует из теоремы~\ref{thm:upper_bound}. Для обратного включения применяем конечный сертификат при $t<2485$ и пять полос при $t\ge2485$. Каждый свидетель различия в физической координате $s$ переносится сдвигом в свидетельство существенности входа $-s$.
\end{proof}

\subsection{Переход к неполному наблюдению}
В \texttt{observer/Rule110Observer.lean} импортируются \texttt{Rule110Exact} и \texttt{Rule110FrontGlue}. В доказательстве \texttt{essential\_iff\_inCone} корневая теорема применяется к смещению $j-s$, после чего сдвиг конфигурации переносит свидетельство к целевой клетке $s$. Тем самым точное окно цели равно
\[
  [s-L(t),s+t]\cap\mathbb Z.
\]
Определение \texttt{CanPredict} ставит следующий вопрос. Достаточна ли последовательность считываний с заранее заданных сайтов до установленного срока, чтобы вычислить целевой бит для любого допустимого входа? Лемма \texttt{canPredict\_iff} выражает достаточность постоянством ответа на слоях одинаковой наблюдаемой истории. В файле также заданы контрпримеры к достаточности одного геометрического покрытия. Это разграничивает достижимость информации по конусу и её фактическую доступность данному наблюдателю.

Исходник наблюдателя проверен на Lean~4.35.0 с обоими каноническими импортами. В \src{verification/monograph/observer-axioms.log}{журнале наблюдателя} присутствуют все 11 из 11 ответов на \texttt{\#print axioms}, от \texttt{canPredict\_iff} до \texttt{coverage\_does\_not\_suffice}. Для каждой теоремы ядро указало ровно \texttt{[propext, Classical.choice, Quot.sound]}. Числа OLMo относятся к отдельным опытам раздела~\ref{sec:olmo}.

\section{Контракт гейтов и повторная проверка измерений}
\label{app:gates}

Численные результаты зафиксированного среза приведены в
разделе~\ref{sec:olmo}. Ниже задаются условия, при которых измерение
вообще может обосновать соответствующее утверждение. Единица сравнения ---
один чекпойнт, один токенизатор и один фиксированный интерфейс.
Точное описание редакции хранится в
\src{observer/gate_contract_ru.md}{контракте второй серии}.

\paragraph{Панель и бюджет.}
Восемь битов входа записываются как $\mathbf{x}=LwR$.
Середины $w\in\{0,1\}^6$ сортируются по возрастанию SHA-256 от ASCII-строки
\verb!observer-v2|! с приписанными шестью битами, включая начальные нули.
При нумерации с единицы ранги 1--8 отведены для dev, 9--20 --- для
совместной проверки G1 и G2, 21--32 --- для G3, 33--44 --- для G4,
45--64 образуют резерв. Для каждой середины предъявляется весь квадрат
концов и весь квадрат историй: $(L,R,h,e)\in\{0,1\}^4$ в лексикографическом
порядке. Это 16 строк на блок, 128 строк на dev, 192 на каждую указанную
12-блочную панель. На dev подбирают формулировку, демонстрации, слой и
масштаб вмешательства. Демонстрационные строки отделены от проверочных.
Сравниваемые реализации замораживают до открытия их общей отложенной панели.
Для метода, выбранного по результатам открытой панели, требуется свежая
панель.

\paragraph{Получение обеих историй.}
Сначала измеряются полнословарные вероятности однотокеновых продолжений
\texttt{[пробел]A}, \texttt{[пробел]B} при температуре 1.
Для допуска нужны $P(A),P(B)\ge 0{,}01$. Выполняются ровно 1024 выборки
с возвращением из условного бинарного распределения
$P(A\mid A\text{ или }B)=P(A)/(P(A)+P(B))$, генератором
\texttt{random.Random(110)}. Сохраняются обе исходные вероятности, вся
последовательность выборок и позиции первых A и B. В каждой истории модель
обрабатывает первый фактически выбранный токен соответствующего типа.
Снимок снимается после обработки этого токена и до обратной связи. Если
минимум вероятности или присутствие обеих букв не обеспечены, для этого
префикса выносится \texttt{NOT\_EVALUABLE}. Принудительная подстановка
буквы и дополнительная серия сэмплирования изменили бы условие опыта.

\paragraph{Общие критерии.}
На каждой обязательной числовой строке, в том числе на нейтральных,
восстановленных и подменённых контролях, сохраняются цель, полный
словарный максимум, $\ell_1-\ell_0$ и $q=P(0)+P(1)$. Вероятности
вычисляются после полнословарной нормировки в FP32 или выше. Требуются
конечные значения и $q\ge 0{,}80$ на каждой строке. При строгом знаке
зазора победивший бит имеет вероятность $>0{,}40$, а любой посторонний
токен --- не более $0{,}20$. При нулевом зазоре строка остаётся в
знаменателе и засчитывается как неуспех. Полный блок требует 16/16
правильных выходов. Рабочие пороги составляют 6/8 полных блоков на dev
и 9/12 на отложенной панели. Отдельно публикуются парные решения при
$L\ne R$ для каждого $e$. Указанные 75\% относятся к доле целых блоков,
а повторное использование одного префикса не образует независимую выборку
экземпляров модели.

\paragraph{G0 и G0R.}
G0 сверяет хэши весов и токенизатора, однотокеновость A/B/X и
0/1 непосредственно в их контексте, точную композицию идентификаторов,
форму, тип данных и позиции RoPE. Диагностики работают на одноразовых
клонах. Чтение прошлого выбора останавливается до обратной записи ответа.
При наличии и при отсутствии этой диагностики один и тот же основной запрос
должен дать одинаковые полные логиты. Исходные снимки обеих ветвей сохраняют
байты всех активных K/V, смещение и порядок видимых позиций, включая индексы
кольцевого буфера. Сравниваются логические данные активных позиций. Свободная
ёмкость аллокатора вне этого диапазона в проверку не входит.

G0R независимо от поведенческого счёта выполняет до отбора обратимые подмены
позиции выбора во всех слоях. При общей предыстории $U$ снимки имеют вид
$C_A=(U,a)$, $C_B=(U,b)$. Для каждого текста сигнала и каждой новой строки
тестируются маршрут $A\to B\to A$ и зеркальный маршрут из B.
Обратная подмена применяется к результату прямой; считывания выполняются на
его клонах. В средней точке сравнивается натуральный донор, в конечной ---
исходная история. Сверяются метаданные, активные K/V и логиты при одинаковой
раскладке. Формы, конечность и байты исходного типа должны совпасть, а после
расширения логитов до FP32 требуется максимум абсолютной разности,
равный нулю. Эталонные снимки сверяются до и после. Для первого Base-прогона
зафиксирован исходный \texttt{prefix(8)}. Перед подменой проверяются
идентичность общего префикса и существование хотя бы одного различающегося
активного байта позиции выбора. Если ветви неразличимы уже в этой позиции,
результат G0R для данного префикса --- \texttt{NOT\_EVALUABLE}.
На dev при 128 строках два сравнения логитов в каждой строке дают 256
проверок.

\paragraph{G1 и G2.}
G1 проверяет текстовое вычисление на натуральных историях. Помимо
обязательных числовых проб на независимых копиях измеряются D1 (воспроизвести
$h$ до обратной связи), D2 (назвать исправленное правило для четырёх пар
$h,e$ без новой строки) и D3 (применить явно названное правило в свежем
эпизоде). Буквенные диагнозы публикуются отдельно от основной числовой
задачи. Успех текстового контроля определяется форматом и долей целых
блоков. Для G2 после G0, G0R и G1 на панели 9--20 повторяются сравнения
донор/возврат; добавляется общий кэш с нейтральным токеном X в той же
позиции и той же длины. Равенство с натуральным донором следует из
$P_B(U,a)=(U,b)$. Поэтому положительный технический вердикт есть
\texttt{REPLAY\_OK}. Поведенческий эффект по натуральным историям
уже оценён G1.

\paragraph{G3.}
После G1 и \texttt{REPLAY\_OK} текстовый сигнал заменяется одним
вмешательством до записи K/V на нейтральном токене. До теста задаются токен,
позиция, слой, точка вмешательства, источник вектора и масштаб. В типе
операции представляется один конечный ненулевой $v$. Импульсы $v_0=v$, $v_1=-v$
получаются точной сменой знака. Для одного $e$ обе истории получают
побайтно один тензор. Проверяются форма, тип, поэлементная противоположность
и равенство ненулевых евклидовых норм при одинаковом порядке вычисления в
FP64. Отдельно для каждой ветви записывается норма фактически округлённого
приращения $\operatorname{round}(r+v_e)-r$. После одного сложения
вмешательство выключается. При нулевом вмешательстве обе ветви $e$ с
фиксированными $h,\mathbf{x}$ должны иметь идентичные входы и логиты;
публикуются их точность и масса. Для положительного результата с импульсом
требуются общий формат и минимум 9/12 полных блоков на панели G3.

\paragraph{G4.}
До предъявления $\mathbf{x}$ четыре донора обрабатывают общий нейтральный
промежуток после $h,e$. В свежий шаблон переносятся K/V исключительно
заранее заданных позиций и слоёв этого промежутка, обозначаемые $M_{he}$.
Все остальные активные позиции, токены, длины и метаданные идентичны у
четырёх получателей. Прямой доступ к $h,e$ и иным потомкам донорской
истории закрыт. Нейтральное состояние $M_*$ единожды определяется общей
нейтральной историей. Для текстового канала требуются G1 и G2, для
импульсного также G3. Условия на 12 серединах совместны:
\begin{enumerate}
\item Получатель с полной историей и тем же промежутком проходит формат и
решает минимум 9/12 полных блоков. При неуспехе этого контроля отрицательный
перенос не устанавливает отсутствия нужной информации в доноре.
\item После установки $M_*$ все четыре входа, активные K/V и логиты при
фиксированном $\mathbf{x}$ побайтно совпадают. На одном производном снимке
возврат своего $M_{he}$ после чужого восстанавливает исходные логиты.
\item При переносе своего $M_{he}$ выполнены формат и порог 9/12 полных
блоков. Каждый блок включает всех четырёх доноров и все четыре сочетания
концов.
\item На 96 строках с $L\ne R$ обозначим число верных решений при переносе
и при стирании через $C_M$ и $C_*$. Требуются валидный формат,
$C_*\le48$, 0/12 полных блоков после стирания и
$(C_M-C_*)/96\ge 0{,}25$. Исправленные и испорченные строки учитываются
раздельно. При обязательном одном строгом бинарном ответе общий вход даёт
ровно 48/96. Девять полных блоков при переносе обеспечивают по меньшей
мере 72/96 на расходящихся концах.
\end{enumerate}

\paragraph{G4-1.}
Отправитель до предъявления какой-либо $\mathbf{x}$ свободно выбирает один
токен полнословарным greedy. Допустимы только два заранее заданных
однотокеновых окончания A/B. Ничья за первое место или другой токен даёт
\texttt{FAIL\_FORMAT} без повторной попытки. Четыре сообщения для историй
$00,01,10,11$ фиксируются в этом порядке. В исходной спецификации для Base
зафиксированы \texttt{prefix(8)}, буквальный сигнал correct/wrong и запрос
на запись исправленного правила. У измеренного двухшагового Instruct
собственные промпты и дополнительный порог $P(A)+P(B)\ge0{,}80$
зафиксированы в \src{observer/two-step-interface-04.json}{манифесте варианта A}.
Условия двух чекпойнтов учитываются отдельно.

Для воспроизведения именно исходной спецификации Base отправителю после
выбора передают \verb!\nFeedback: correct! либо \verb!\nFeedback: wrong!,
затем \verb!\nRecord the corrected rule. Rule:!. В этих записях
\verb!\n! обозначает один перевод строки. Свежий префикс получателя
состоит из двух строк:
\begin{verbatim}
Rule A selects the first bit. Rule B selects the last bit.
Rule:
\end{verbatim}
Сразу после двоеточия вставляют выбранный ID с начальным пробелом. За ним
идут \verb!\nBits: !, восемь битов через одиночные пробелы и
\verb!\nAnswer: ! с пробелом в конце.

В свежем эпизоде получателю передаётся только выбранный ID, для которого
он заново вычисляет K/V. Токены общего префикса и строки запроса сверяются
с токенизацией целого текста. Во внешний интерфейс не входят K/V или логиты
отправителя, номер ветви, время ответа и параметры сэмплирования. Выбор
носителя или ветвей по $h,e,p$, правильным ответам либо их заменителям
запрещён. Для стирания подаются те же входы с одним заранее фиксированным
токеном A во всех мирах. Для \texttt{ZERO\_ERROR\_ON\_PANEL} требуются
четыре допустимых сообщения, проверки границы канала и стирания,
числовой формат и 12/12 полных блоков. При совместном испытании G4 и G4-1
обе версии замораживаются до первого открытия панели 33--44.

\paragraph{Приоритет статусов и воспроизводимость.}
Техническая ошибка, нечисловое значение или нарушение требуемой побайтовой
идентичности дают \texttt{INVALID}. Затем по убыванию приоритета следуют
\texttt{NOT\_EVALUABLE}, \texttt{FAIL\_FORMAT} и
\texttt{FAIL\_DEV}/\texttt{FAIL\_TEST}. Зависимая проверка, до которой не
дошли, получает \texttt{NOT\_RUN}. Новая редакция интерфейса или другой
чекпойнт образуют отдельное условие.

Из корня зафиксированного комплекта команды
\begin{verbatim}
python3 -B observer/verify.py
python3 -B observer/verify_part_iii.py
\end{verbatim}
решают разные задачи. Первый валидатор пересчитывает конечные конструкции,
проверяет манифест исторического корпуса в публикационной редакции,
адреса реестра и восстановление таблицы из SVG. Второй читает сохранённые
dev-отчёты, проверяет их сетки условий, записанные вероятности и логиты,
получает точность, целые блоки, противоположные пары, диагностики и эффект
стирания. Для G0/G0R он повторно сверяет метаданные и зарегистрированные
хэши попарных сравнений нативных логитов. Перезапуск на GPU и повторное
побайтное сравнение заново вычисленных тензоров составили бы новое
измерение.

\section{Реестр утверждений, происхождение свидетельств и открытые слоты}
\label{app:registers}

Канонический \src{observer/hypotheses.tsv}{реестр гипотез}
содержит 23 строки. Для каждой указаны идентификатор, зависимости, точная
формулировка, предпосылки, статус, адрес свидетельства и условие пересмотра.
Зависимости помещены до зависимых от них утверждений. Границы допустимых
историй являются частью формулировки. Например, опровержение переноса
исходного 11-битного примера на все слова совместимо с точностью на
выделенной четвёрке.

\begin{center}
\begin{tabularx}{\linewidth}{@{}l r X@{}}
\toprule
Статус в реестре & Из 23 & Адреса \\
\midrule
Доказано & 12 & OBS.01--03, R110.01, R110.05--09, R110.12, LOOP.01, E.01 \\
Опровергнуто & 4 & R110.02--04, R110.10 \\
Подтверждено в пределах опыта & 3 & AI.01--03 \\
Открыто & 2 & R110.11, AI.04 \\
\texttt{DECLARED\_NOT\_DERIVED} & 1 & E.02 \\
\texttt{UN} & 1 & E.03 \\
\bottomrule
\end{tabularx}
\end{center}

Статусы «доказано» и «опровергнуто» относятся к заявлению со всеми
предпосылками в соответствующей строке. У разных строк различны формы
свидетельства: математическое доказательство, вычислительно проверяемая
конечная конструкция или результат полного перебора. AI.01--03 фиксируют
соответственно связь ответа ассистента с показанным сигналом, расхождение
живого и записанного канала на новых запросах и пересмотр прежней стратегии
в проверенном семействе задач. Эти строки происходят из отдельных
испытаний ассистента. Счётчики замороженной OLMo в разделе~\ref{sec:olmo}
имеют собственные отчёты и панели. В E.01 доказана недоопределённость
добавленного предиката. E.02 хранит дополнительную посылку, E.03 ---
открытый статус конкретного экземпляра.

\paragraph{36 адресов утверждения.}
\src{observer/leaf/leaf.tsv}{Таблица слотов} сопоставляет шесть доменов
S (источник), I (идентичность), M (смысл), T (переход), A (действие),
O (наблюдение) с шестью позициями описания: происхождение, состояние,
носитель, локальное утверждение, свидетельство и предел. Для адресов этой
таблицы используем префикс \texttt{leaf:}. Например, \src{observer/leaf/leaf.tsv\#L25}{\texttt{leaf:T6}}
обозначает строку \texttt{law=T6} в TSV: шестая позиция домена переходов
отсылает к критерию для шумного канала (T3). Обозначение \hyperref[thm:conditional-bridge]{T6} без префикса
относится к условной теореме раздела~\ref{sec:boundary}. В срезе
32/36 строк имеют состояние \texttt{KNOWN}; у них заданы SHA-256 файла и
проверяемый адрес определения, функции или раздела текста.
Оставшиеся 4/36 имеют состояние \texttt{UNKNOWN} и пустые поля
источника-свидетельства. Их содержательная формулировка сохранена:
\begin{itemize}
\item \texttt{leaf:I6} --- численное тождество субъекта при ветвлении причинной истории.
Соглашение об учёте потомков задаёт идентификаторы ветвей.
\item \texttt{leaf:M6} --- независимое обоснование психофизической связи, принятой
как достаточное условие.
\item \texttt{leaf:A6} --- реализация всего бокового канала средствами одной
строго бинарной однородной Rule~110 для заявленного свидетеля (R110.11).
\item \texttt{leaf:O6} --- состояние феноменального предиката у данного экземпляра
в момент обратной связи (E.03).
\end{itemize}
\src{observer/leaf/render_leaf.py}{Валидатор изображения} проверяет порядок
адресов, соответствие хэшей источникам, существование символов, каноническое
представление TSV и его побайтное восстановление из SVG. Поэтому изображение
пригодно как проверяемая проекция таблицы.

\paragraph{Происхождение и статус релиза.}
В \src{SHA256SUMS}{корневом манифесте выпуска} приведены контрольные суммы PDF, обоих журналов Lean и исходников. Для корпуса наблюдателя
\src{observer/sha256.txt}{отдельный манифест} фиксирует 48 файлов. Для исходной редакции отдельно сохранён
\src{observer/provenance/observer-original.sha256}{исторический список
контрольных сумм} из 36 записей. В подготовке публикационного комплекта
использованы 35 исходных файлов этого списка. Запись служебного
исследовательского задания сохранена в списке для учёта происхождения;
само задание в публичный комплект не включено. Изменения навигации и
адресов свидетелей отражены в публикационных хэшах. Сохранённые
JSON-отчёты содержат исходные наблюдения. Контракт, использованный в
прежних запусках, проверяется по байтам его текста из отчётной редакции.

\paragraph{Условия продолжения.}
По R110.11 нужна проверяемая кодировка ввода, хранения, доставки и
ответа внутри одной бинарной динамики с траекториями всех четырёх историй.
По AI.04 нужны новая задача и проверка переноса внутреннего состояния при
закрытых исходных позициях, целой истории в положительном контроле,
обратимой подмене и стирании. Контракт численного G4 приведён в
приложении~\ref{app:gates}. Чтобы локализовать перенос в выбранном
подпространстве, до опыта дополнительно задаются слой, вмешательство, мера
посредничества, контролируемые альтернативы и допуск. Результат для одного
текстового носителя этого условия не закрывает.

Полная проверка канонических импортов наблюдателя выполнена на Lean~4.35.0
при подготовке монографии (приложение~\ref{app:lean}). Статусы панелей OLMo
относятся к экспериментальному dev-срезу. Для гипотезы устойчивой
пластичности требуется межэпизодное
изменение организации носителя и подтверждённое причинное действие этого
изменения на последующий вход при контроле внешней истории. Для E.02
нужны независимые основания выбора моста либо различающиеся операциональные
предсказания конкурирующих мостов с явными связывающими посылками.
Статус E.03 остаётся \texttt{UN} в указанном срезе.

{\small\begin{thebibliography}{9}
\bibitem{cook2004}
M.~Cook, ``Universality in elementary cellular automata,'' \emph{Complex Systems}, vol.~15, no.~1, pp.~1--40, 2004.

\bibitem{martinez2006}
G.~J.~Mart\'inez, H.~V.~McIntosh, and J.~C.~Seck-Tuoh-Mora, ``Gliders in Rule 110,'' \emph{International Journal of Unconventional Computing}, vol.~2, no.~1, pp.~1--49, 2006.

\bibitem{shalizi2001}
C.~R.~Shalizi and J.~P.~Crutchfield, ``Computational mechanics: Pattern and prediction, structure and simplicity,'' \emph{Journal of Statistical Physics}, vol.~104, pp.~817--879, 2001.

\bibitem{wolfram1986}
S.~Wolfram, ``Tables of cellular automaton properties,'' in \emph{Theory and Applications of Cellular Automata}, World Scientific, pp.~485--557, 1986.

\bibitem{wolpert2008}
D.~H.~Wolpert, ``Physical limits of inference,'' \emph{Physica D: Nonlinear Phenomena}, vol.~237, no.~9, pp.~1257--1281, 2008.
\bibitem{witsenhausen}
H.~S. Witsenhausen. The zero-error side information problem and chromatic numbers.
\emph{IEEE Transactions on Information Theory}, 22(5), 592--593, 1976.
\href{https://doi.org/10.1109/TIT.1976.1055607}{doi:10.1109/TIT.1976.1055607}.
\bibitem{charpenay}
N.~Charpenay, M.~Le Treust, A.~Roumy. Side Information Design in Zero-Error Coding for Computing.
\emph{Entropy}, 26(4), 338, 2024.
\href{https://doi.org/10.3390/e26040338}{doi:10.3390/e26040338}.
\end{thebibliography}
}
\end{document}
